%%%%%%%%%%%%%%%%%%%%%%%%%%%%%%%%%%%%%%%%%%%%%%%%%%%%%%%%%%%%
% 3.7 DIFFERENTIAL GEOMETRY
% The Composition of Advanced Mathematics
%%%%%%%%%%%%%%%%%%%%%%%%%%%%%%%%%%%%%%%%%%%%%%%%%%%%%%%%%%%%

\section{Differential Geometry}
\label{sec:differential-geometry}

\prerequisite{
Multivariable Calculus from \cref{sec:multivariable-calculus},
Vector Calculus from \cref{sec:vector-calculus},
Linear Algebra from \cref{sec:linear-algebra},
Euclidean Geometry from \cref{sec:euclidean-geometry},
and Non-Euclidean Geometry from
\cref{sec:non-euclidean-geometry}.
}

\learningobjectives{
By the end of this section, the reader should be able to:
\begin{itemize}
    \item explain the goals of differential geometry;
    \item describe smooth curves and regular parametrizations;
    \item compute tangent, normal, and binormal vectors;
    \item calculate arc length and curvature of curves;
    \item interpret torsion and the Frenet--Serret frame;
    \item describe parametrized surfaces and tangent planes;
    \item compute surface normals and surface area;
    \item distinguish the first and second fundamental forms;
    \item calculate Gaussian and mean curvature in basic examples;
    \item distinguish intrinsic and extrinsic curvature;
    \item explain principal curvatures and principal directions;
    \item describe geodesics and geodesic distance;
    \item recognize manifolds, coordinate charts, and tangent spaces;
    \item understand the role of Riemannian metrics;
    \item explain covariant differentiation and parallel transport;
    \item distinguish Gaussian, sectional, Ricci, and scalar curvature;
    \item state the geometric meaning of the Gauss--Bonnet theorem;
    \item connect differential geometry with topology, relativity,
    optimization, mechanics, and geometric analysis.
\end{itemize}
}

%%%%%%%%%%%%%%%%%%%%%%%%%%%%%%%%%%%%%%%%%%%%%%%%%%%%%%%%%%%%
\subsection{What Is Differential Geometry?}
\label{subsec:what-is-differential-geometry}
%%%%%%%%%%%%%%%%%%%%%%%%%%%%%%%%%%%%%%%%%%%%%%%%%%%%%%%%%%%%

Differential geometry combines geometry with the methods of calculus and
linear algebra.

Instead of studying only rigid geometric figures, differential geometry
studies smooth curves, surfaces, and higher-dimensional spaces whose geometry
may vary continuously from point to point.

\begin{definitionbox}{Differential Geometry}
\emph{Differential geometry} is the study of geometric objects using
derivatives, integrals, differential equations, linear algebra, and related
analytic methods.
\end{definitionbox}

Some of its central questions are:

\begin{itemize}
    \item How does a curve bend?
    \item How does a surface bend?
    \item What is the shortest path between two points?
    \item How can curvature be measured intrinsically?
    \item How can calculus be performed on curved spaces?
    \item How does local curvature influence global topology?
\end{itemize}

%%%%%%%%%%%%%%%%%%%%%%%%%%%%%%%%%%%%%%%%%%%%%%%%%%%%%%%%%%%%
\subsection{From Calculus to Geometry}
\label{subsec:calculus-to-geometry}
%%%%%%%%%%%%%%%%%%%%%%%%%%%%%%%%%%%%%%%%%%%%%%%%%%%%%%%%%%%%

Ordinary calculus studies functions such as

\[
f:\R\to\R
\]

and

\[
f:\R^n\to\R^m.
\]

Differential geometry uses these same tools to study geometric objects.

A parametrized curve is a map

\[
\gamma:I\to\R^n,
\]

while a parametrized surface is locally described by a map

\[
\mathbf{X}:U\subseteq\R^2\to\R^3.
\]

More generally, a smooth \(n\)-dimensional manifold locally resembles

\[
\R^n.
\]

Thus calculus becomes a local language for studying spaces that may be
globally curved.

%%%%%%%%%%%%%%%%%%%%%%%%%%%%%%%%%%%%%%%%%%%%%%%%%%%%%%%%%%%%
\subsection{Smooth Curves}
\label{subsec:smooth-curves}
%%%%%%%%%%%%%%%%%%%%%%%%%%%%%%%%%%%%%%%%%%%%%%%%%%%%%%%%%%%%

\begin{definitionbox}{Parametrized Curve}
A \emph{parametrized curve} in \(\R^n\) is a map

\[
\boxed{
\gamma:I\to\R^n,
}
\]

where \(I\subseteq\R\) is an interval.
\end{definitionbox}

The Greek letter

\[
\gamma
\]

is read ``gamma.''

In \(\R^3\), we may write

\[
\boxed{
\gamma(t)
=
\bigl(x(t),y(t),z(t)\bigr).
}
\]

The parameter \(t\) need not represent time, although that interpretation is
common in applications.

%%%%%%%%%%%%%%%%%%%%%%%%%%%%%%%%%%%%%%%%%%%%%%%%%%%%%%%%%%%%
\subsection{Velocity and Tangent Vectors}
\label{subsec:curve-velocity-tangent}
%%%%%%%%%%%%%%%%%%%%%%%%%%%%%%%%%%%%%%%%%%%%%%%%%%%%%%%%%%%%

Differentiating a curve gives

\[
\boxed{
\gamma'(t)
=
\bigl(x'(t),y'(t),z'(t)\bigr).
}
\]

The vector

\[
\gamma'(t)
\]

is tangent to the curve and may also be interpreted as velocity when \(t\)
represents time.

Its magnitude

\[
\boxed{
\|\gamma'(t)\|
}
\]

is the speed.

%%%%%%%%%%%%%%%%%%%%%%%%%%%%%%%%%%%%%%%%%%%%%%%%%%%%%%%%%%%%
\subsection{Regular Curves}
\label{subsec:regular-curves}
%%%%%%%%%%%%%%%%%%%%%%%%%%%%%%%%%%%%%%%%%%%%%%%%%%%%%%%%%%%%

\begin{definitionbox}{Regular Curve}
A smooth parametrized curve is \emph{regular} if

\[
\boxed{
\gamma'(t)\neq\mathbf{0}
}
\]

at every parameter value under consideration.
\end{definitionbox}

Regularity guarantees that the curve has a well-defined tangent direction.

%%%%%%%%%%%%%%%%%%%%%%%%%%%%%%%%%%%%%%%%%%%%%%%%%%%%%%%%%%%%
\subsection{Example: The Helix}
\label{subsec:helix-example}
%%%%%%%%%%%%%%%%%%%%%%%%%%%%%%%%%%%%%%%%%%%%%%%%%%%%%%%%%%%%

Consider

\[
\boxed{
\gamma(t)
=
(a\cos t,a\sin t,bt).
}
\]

This curve is a circular helix.

Its derivative is

\[
\gamma'(t)
=
(-a\sin t,a\cos t,b).
\]

Therefore

\[
\|\gamma'(t)\|
=
\sqrt{a^2\sin^2t+a^2\cos^2t+b^2}
=
\sqrt{a^2+b^2}.
\]

Hence the helix has constant speed.

%%%%%%%%%%%%%%%%%%%%%%%%%%%%%%%%%%%%%%%%%%%%%%%%%%%%%%%%%%%%
\subsection{Arc Length}
\label{subsec:differential-geometry-arc-length}
%%%%%%%%%%%%%%%%%%%%%%%%%%%%%%%%%%%%%%%%%%%%%%%%%%%%%%%%%%%%

The length of a smooth curve

\[
\gamma:[a,b]\to\R^n
\]

is

\[
\boxed{
L
=
\int_a^b
\|\gamma'(t)\|\,dt.
}
\]

This quantity is independent of any orientation-preserving regular
reparametrization.

Thus arc length belongs to the geometric curve itself rather than to a
particular choice of parameter.

%%%%%%%%%%%%%%%%%%%%%%%%%%%%%%%%%%%%%%%%%%%%%%%%%%%%%%%%%%%%
\subsection{Arc-Length Parameter}
\label{subsec:arc-length-parameter}
%%%%%%%%%%%%%%%%%%%%%%%%%%%%%%%%%%%%%%%%%%%%%%%%%%%%%%%%%%%%

Define

\[
\boxed{
s(t)
=
\int_{t_0}^{t}
\|\gamma'(u)\|\,du.
}
\]

The variable \(s\) measures distance traveled along the curve.

When a curve is parametrized directly by arc length,

\[
\boxed{
\left\|
\frac{d\gamma}{ds}
\right\|
=1.
}
\]

Such a parametrization is called a \emph{unit-speed parametrization}.

%%%%%%%%%%%%%%%%%%%%%%%%%%%%%%%%%%%%%%%%%%%%%%%%%%%%%%%%%%%%
\subsection{Unit Tangent Vector}
\label{subsec:unit-tangent-vector}
%%%%%%%%%%%%%%%%%%%%%%%%%%%%%%%%%%%%%%%%%%%%%%%%%%%%%%%%%%%%

For a regular curve, the unit tangent vector is

\[
\boxed{
\mathbf{T}(t)
=
\frac{\gamma'(t)}
{\|\gamma'(t)\|}.
}
\]

The symbol

\[
\mathbf{T}
\]

is read ``T'' and denotes the unit tangent direction.

If the curve is parametrized by arc length,

\[
\boxed{
\mathbf{T}(s)=\gamma'(s).
}
\]

%%%%%%%%%%%%%%%%%%%%%%%%%%%%%%%%%%%%%%%%%%%%%%%%%%%%%%%%%%%%
\subsection{Curvature of a Curve}
\label{subsec:curve-curvature}
%%%%%%%%%%%%%%%%%%%%%%%%%%%%%%%%%%%%%%%%%%%%%%%%%%%%%%%%%%%%

A straight line has constant tangent direction.

A curved path has a tangent direction that changes.

This motivates curvature.

\begin{definitionbox}{Curvature}
For a unit-speed curve, the \emph{curvature} is

\[
\boxed{
\kappa(s)
=
\left\|
\frac{d\mathbf{T}}{ds}
\right\|.
}
\]
\end{definitionbox}

The Greek letter

\[
\kappa
\]

is read ``kappa.''

Large curvature means the tangent direction changes rapidly with arc length.

Small curvature means the curve bends gradually.

%%%%%%%%%%%%%%%%%%%%%%%%%%%%%%%%%%%%%%%%%%%%%%%%%%%%%%%%%%%%
\subsection{Curvature for a General Parametrization}
\label{subsec:general-curve-curvature}
%%%%%%%%%%%%%%%%%%%%%%%%%%%%%%%%%%%%%%%%%%%%%%%%%%%%%%%%%%%%

For a regular space curve

\[
\gamma(t)
\]

in \(\R^3\),

\[
\boxed{
\kappa(t)
=
\frac{
\|\gamma'(t)\times\gamma''(t)\|
}{
\|\gamma'(t)\|^3
}.
}
\]

For a plane curve

\[
\gamma(t)=(x(t),y(t)),
\]

this becomes

\[
\boxed{
\kappa(t)
=
\frac{
|x'(t)y''(t)-y'(t)x''(t)|
}{
\left(x'(t)^2+y'(t)^2\right)^{3/2}
}.
}
\]

%%%%%%%%%%%%%%%%%%%%%%%%%%%%%%%%%%%%%%%%%%%%%%%%%%%%%%%%%%%%
\subsection{Curvature of a Graph}
\label{subsec:graph-curvature}
%%%%%%%%%%%%%%%%%%%%%%%%%%%%%%%%%%%%%%%%%%%%%%%%%%%%%%%%%%%%

If a plane curve is given as

\[
y=f(x),
\]

then

\[
\boxed{
\kappa
=
\frac{|f''(x)|}
{\left(1+[f'(x)]^2\right)^{3/2}}.
}
\]

This formula connects the second derivative with geometric bending.

However, curvature is not simply

\[
f''(x).
\]

The denominator corrects for the slope of the graph.

%%%%%%%%%%%%%%%%%%%%%%%%%%%%%%%%%%%%%%%%%%%%%%%%%%%%%%%%%%%%
\subsection{Worked Example: Curvature of a Parabola}
\label{subsec:worked-curvature-parabola}
%%%%%%%%%%%%%%%%%%%%%%%%%%%%%%%%%%%%%%%%%%%%%%%%%%%%%%%%%%%%

\begin{workedexamplebox}{Curvature of \(y=x^2\)}

For

\[
f(x)=x^2,
\]

we have

\[
f'(x)=2x,
\qquad
f''(x)=2.
\]

Therefore

\[
\kappa(x)
=
\frac{2}
{(1+4x^2)^{3/2}}.
\]

At

\[
x=0,
\]

we obtain

\[
\kappa(0)=2.
\]

\begin{answerbox}
\[
\boxed{\kappa(0)=2}
\]
\end{answerbox}
\end{workedexamplebox}

%%%%%%%%%%%%%%%%%%%%%%%%%%%%%%%%%%%%%%%%%%%%%%%%%%%%%%%%%%%%
\subsection{Radius of Curvature}
\label{subsec:radius-of-curvature}
%%%%%%%%%%%%%%%%%%%%%%%%%%%%%%%%%%%%%%%%%%%%%%%%%%%%%%%%%%%%

When

\[
\kappa\neq0,
\]

the \emph{radius of curvature} is

\[
\boxed{
\rho=\frac{1}{\kappa}.
}
\]

The Greek letter

\[
\rho
\]

is read ``rho.''

A curve with large curvature has a small radius of curvature.

A curve with small curvature has a large radius of curvature.

%%%%%%%%%%%%%%%%%%%%%%%%%%%%%%%%%%%%%%%%%%%%%%%%%%%%%%%%%%%%
\subsection{Normal Vector}
\label{subsec:principal-normal-vector}
%%%%%%%%%%%%%%%%%%%%%%%%%%%%%%%%%%%%%%%%%%%%%%%%%%%%%%%%%%%%

When

\[
\kappa\neq0,
\]

the principal unit normal vector is

\[
\boxed{
\mathbf{N}
=
\frac{\mathbf{T}'}
{\|\mathbf{T}'\|}
}
\]

when differentiation is taken with respect to arc length.

The vector

\[
\mathbf{N}
\]

points in the direction in which the tangent vector is turning.

%%%%%%%%%%%%%%%%%%%%%%%%%%%%%%%%%%%%%%%%%%%%%%%%%%%%%%%%%%%%
\subsection{Binormal Vector}
\label{subsec:binormal-vector}
%%%%%%%%%%%%%%%%%%%%%%%%%%%%%%%%%%%%%%%%%%%%%%%%%%%%%%%%%%%%

For a regular space curve with nonzero curvature, define

\[
\boxed{
\mathbf{B}
=
\mathbf{T}\times\mathbf{N}.
}
\]

The vectors

\[
\mathbf{T},
\qquad
\mathbf{N},
\qquad
\mathbf{B}
\]

form an orthonormal frame along the curve.

This moving coordinate system is called the \emph{Frenet frame}.

%%%%%%%%%%%%%%%%%%%%%%%%%%%%%%%%%%%%%%%%%%%%%%%%%%%%%%%%%%%%
\subsection{The Frenet Frame}
\label{subsec:frenet-frame}
%%%%%%%%%%%%%%%%%%%%%%%%%%%%%%%%%%%%%%%%%%%%%%%%%%%%%%%%%%%%

\begin{centereddiagram}
\begin{tikzpicture}[scale=1.0]
\draw[thick,domain=-2:2,samples=80]
plot (\x,{0.22*\x*\x});

\coordinate (P) at (0,0);
\fill (P) circle (2pt);

\draw[-{Latex[length=3mm]},thick]
(P)--(2,0) node[right] {\(\mathbf{T}\)};

\draw[-{Latex[length=3mm]},thick]
(P)--(0,1.5) node[above] {\(\mathbf{N}\)};

\draw[-{Latex[length=3mm]},thick]
(P)--(-0.75,-0.75) node[left] {\(\mathbf{B}\)};

\node[below right] at (P) {\(\gamma(s)\)};
\end{tikzpicture}
\end{centereddiagram}

The diagram is schematic.

For a genuine space curve, the binormal is perpendicular to the osculating
plane spanned by \(\mathbf{T}\) and \(\mathbf{N}\).

%%%%%%%%%%%%%%%%%%%%%%%%%%%%%%%%%%%%%%%%%%%%%%%%%%%%%%%%%%%%
\subsection{Torsion}
\label{subsec:curve-torsion}
%%%%%%%%%%%%%%%%%%%%%%%%%%%%%%%%%%%%%%%%%%%%%%%%%%%%%%%%%%%%

Curvature measures bending.

Torsion measures twisting out of the osculating plane.

\begin{definitionbox}{Torsion}
For an arc-length parametrized curve with nonzero curvature, the
\emph{torsion} \(\tau\) is defined by

\[
\boxed{
\tau
=
-
\frac{d\mathbf{B}}{ds}\cdot\mathbf{N}.
}
\]
\end{definitionbox}

The Greek letter

\[
\tau
\]

is read ``tau.''

For a sufficiently regular general parametrization,

\[
\boxed{
\tau(t)
=
\frac{
\det\bigl(\gamma'(t),\gamma''(t),\gamma'''(t)\bigr)
}{
\|\gamma'(t)\times\gamma''(t)\|^2
}.
}
\]

%%%%%%%%%%%%%%%%%%%%%%%%%%%%%%%%%%%%%%%%%%%%%%%%%%%%%%%%%%%%
\subsection{Frenet--Serret Formulas}
\label{subsec:frenet-serret}
%%%%%%%%%%%%%%%%%%%%%%%%%%%%%%%%%%%%%%%%%%%%%%%%%%%%%%%%%%%%

For a unit-speed curve with positive curvature,

\[
\boxed{
\frac{d\mathbf{T}}{ds}
=
\kappa\mathbf{N},
}
\]

\[
\boxed{
\frac{d\mathbf{N}}{ds}
=
-\kappa\mathbf{T}
+
\tau\mathbf{B},
}
\]

and

\[
\boxed{
\frac{d\mathbf{B}}{ds}
=
-\tau\mathbf{N}.
}
\]

These are the \emph{Frenet--Serret equations}.

They describe how the moving frame changes along the curve.

%%%%%%%%%%%%%%%%%%%%%%%%%%%%%%%%%%%%%%%%%%%%%%%%%%%%%%%%%%%%
\subsection{Worked Example: Helix Curvature and Torsion}
\label{subsec:worked-helix-curvature-torsion}
%%%%%%%%%%%%%%%%%%%%%%%%%%%%%%%%%%%%%%%%%%%%%%%%%%%%%%%%%%%%

\begin{workedexamplebox}{Curvature and Torsion of a Helix}

Consider

\[
\gamma(t)
=
(a\cos t,a\sin t,bt).
\]

Then

\[
\gamma'(t)
=
(-a\sin t,a\cos t,b),
\]

\[
\gamma''(t)
=
(-a\cos t,-a\sin t,0),
\]

and

\[
\gamma'''(t)
=
(a\sin t,-a\cos t,0).
\]

The resulting curvature and torsion are

\[
\kappa
=
\frac{|a|}
{a^2+b^2},
\]

and, with the orientation determined by the given parametrization,

\[
\tau
=
\frac{b}
{a^2+b^2}.
\]

\begin{answerbox}
\[
\boxed{
\kappa=\frac{|a|}{a^2+b^2},
\qquad
\tau=\frac{b}{a^2+b^2}
}
\]
\end{answerbox}
\end{workedexamplebox}

%%%%%%%%%%%%%%%%%%%%%%%%%%%%%%%%%%%%%%%%%%%%%%%%%%%%%%%%%%%%
\subsection{The Fundamental Theorem of Space Curves}
\label{subsec:fundamental-theorem-space-curves}
%%%%%%%%%%%%%%%%%%%%%%%%%%%%%%%%%%%%%%%%%%%%%%%%%%%%%%%%%%%%

Curvature and torsion contain essentially all local geometric information
about a sufficiently regular space curve.

\begin{theorem}[Fundamental Theorem of Space Curves]
Given suitable functions

\[
\kappa(s)>0
\]

and

\[
\tau(s),
\]

there exists, locally, a unit-speed space curve having curvature
\(\kappa\) and torsion \(\tau\).

The curve is unique up to an orientation-preserving Euclidean rigid motion.
\end{theorem}

Thus

\[
\boxed{
\kappa(s)
\text{ and }
\tau(s)
}
\]

encode the local shape of a space curve.

%%%%%%%%%%%%%%%%%%%%%%%%%%%%%%%%%%%%%%%%%%%%%%%%%%%%%%%%%%%%
\subsection{From Curves to Surfaces}
\label{subsec:curves-to-surfaces}
%%%%%%%%%%%%%%%%%%%%%%%%%%%%%%%%%%%%%%%%%%%%%%%%%%%%%%%%%%%%

A surface requires two independent parameters.

A parametrized surface in \(\R^3\) has the form

\[
\boxed{
\mathbf{X}(u,v)
=
\bigl(
x(u,v),
y(u,v),
z(u,v)
\bigr).
}
\]

The variables

\[
u,v
\]

serve as local coordinates on the surface.

%%%%%%%%%%%%%%%%%%%%%%%%%%%%%%%%%%%%%%%%%%%%%%%%%%%%%%%%%%%%
\subsection{Regular Parametrized Surfaces}
\label{subsec:regular-parametrized-surfaces}
%%%%%%%%%%%%%%%%%%%%%%%%%%%%%%%%%%%%%%%%%%%%%%%%%%%%%%%%%%%%

Differentiate with respect to the two coordinates:

\[
\mathbf{X}_u
=
\frac{\partial\mathbf{X}}{\partial u},
\]

\[
\mathbf{X}_v
=
\frac{\partial\mathbf{X}}{\partial v}.
\]

\begin{definitionbox}{Regular Surface Patch}
A parametrized surface patch is \emph{regular} when

\[
\boxed{
\mathbf{X}_u\times\mathbf{X}_v\neq\mathbf{0}.
}
\]
\end{definitionbox}

This guarantees that the two coordinate directions span a genuine tangent
plane.

%%%%%%%%%%%%%%%%%%%%%%%%%%%%%%%%%%%%%%%%%%%%%%%%%%%%%%%%%%%%
\subsection{Tangent Plane}
\label{subsec:surface-tangent-plane}
%%%%%%%%%%%%%%%%%%%%%%%%%%%%%%%%%%%%%%%%%%%%%%%%%%%%%%%%%%%%

At a regular point

\[
p=\mathbf{X}(u_0,v_0),
\]

the tangent plane is spanned by

\[
\boxed{
\mathbf{X}_u(u_0,v_0)
\qquad\text{and}\qquad
\mathbf{X}_v(u_0,v_0).
}
\]

Thus

\[
T_pS
=
\operatorname{span}
\{
\mathbf{X}_u,\mathbf{X}_v
\}.
\]

The notation

\[
T_pS
\]

is read ``the tangent space to \(S\) at \(p\).''

%%%%%%%%%%%%%%%%%%%%%%%%%%%%%%%%%%%%%%%%%%%%%%%%%%%%%%%%%%%%
\subsection{Surface Normal}
\label{subsec:surface-normal}
%%%%%%%%%%%%%%%%%%%%%%%%%%%%%%%%%%%%%%%%%%%%%%%%%%%%%%%%%%%%

A normal vector is

\[
\boxed{
\mathbf{X}_u\times\mathbf{X}_v.
}
\]

A corresponding unit normal is

\[
\boxed{
\mathbf{n}
=
\frac{
\mathbf{X}_u\times\mathbf{X}_v
}{
\|\mathbf{X}_u\times\mathbf{X}_v\|
}.
}
\]

The choice of

\[
\mathbf{n}
\]

or

\[
-\mathbf{n}
\]

corresponds to choosing an orientation locally.

%%%%%%%%%%%%%%%%%%%%%%%%%%%%%%%%%%%%%%%%%%%%%%%%%%%%%%%%%%%%
\subsection{Example: The Sphere}
\label{subsec:sphere-parametrization}
%%%%%%%%%%%%%%%%%%%%%%%%%%%%%%%%%%%%%%%%%%%%%%%%%%%%%%%%%%%%

A sphere of radius \(R\) may be parametrized by

\[
\boxed{
\mathbf{X}(\theta,\phi)
=
\left(
R\sin\phi\cos\theta,
R\sin\phi\sin\theta,
R\cos\phi
\right).
}
\]

Here

\[
0<\phi<\pi,
\qquad
0<\theta<2\pi
\]

on a standard coordinate patch.

The parameters are angular coordinates.

%%%%%%%%%%%%%%%%%%%%%%%%%%%%%%%%%%%%%%%%%%%%%%%%%%%%%%%%%%%%
\subsection{Surface Area Element}
\label{subsec:surface-area-element}
%%%%%%%%%%%%%%%%%%%%%%%%%%%%%%%%%%%%%%%%%%%%%%%%%%%%%%%%%%%%

For a regular parametrized surface,

\[
\boxed{
dA
=
\|\mathbf{X}_u\times\mathbf{X}_v\|
\,du\,dv.
}
\]

Therefore the area of a parametrized region is

\[
\boxed{
A
=
\iint_D
\|\mathbf{X}_u\times\mathbf{X}_v\|
\,du\,dv.
}
\]

This formula is familiar from multivariable calculus but now has a geometric
interpretation in terms of the tangent vectors of the surface.

%%%%%%%%%%%%%%%%%%%%%%%%%%%%%%%%%%%%%%%%%%%%%%%%%%%%%%%%%%%%
\subsection{The First Fundamental Form}
\label{subsec:first-fundamental-form}
%%%%%%%%%%%%%%%%%%%%%%%%%%%%%%%%%%%%%%%%%%%%%%%%%%%%%%%%%%%%

The tangent vectors determine how lengths and angles are measured on the
surface.

Define

\[
\boxed{
E=\mathbf{X}_u\cdot\mathbf{X}_u,
}
\]

\[
\boxed{
F=\mathbf{X}_u\cdot\mathbf{X}_v,
}
\]

and

\[
\boxed{
G=\mathbf{X}_v\cdot\mathbf{X}_v.
}
\]

Then

\[
\boxed{
ds^2
=
E\,du^2
+
2F\,du\,dv
+
G\,dv^2.
}
\]

This quadratic form is called the \emph{first fundamental form}.

%%%%%%%%%%%%%%%%%%%%%%%%%%%%%%%%%%%%%%%%%%%%%%%%%%%%%%%%%%%%
\subsection{Meaning of the First Fundamental Form}
\label{subsec:first-form-meaning}
%%%%%%%%%%%%%%%%%%%%%%%%%%%%%%%%%%%%%%%%%%%%%%%%%%%%%%%%%%%%

The first fundamental form determines intrinsic metric quantities including:

\begin{itemize}
    \item lengths of tangent vectors;
    \item lengths of curves;
    \item angles;
    \item surface area;
    \item geodesic distance.
\end{itemize}

In matrix form,

\[
\boxed{
g=
\begin{pmatrix}
E&F\\
F&G
\end{pmatrix}.
}
\]

This is the local metric matrix.

%%%%%%%%%%%%%%%%%%%%%%%%%%%%%%%%%%%%%%%%%%%%%%%%%%%%%%%%%%%%
\subsection{Area from the First Fundamental Form}
\label{subsec:area-first-form}
%%%%%%%%%%%%%%%%%%%%%%%%%%%%%%%%%%%%%%%%%%%%%%%%%%%%%%%%%%%%

Since

\[
\|\mathbf{X}_u\times\mathbf{X}_v\|^2
=
EG-F^2,
\]

the area element may be written

\[
\boxed{
dA
=
\sqrt{EG-F^2}\,du\,dv.
}
\]

Thus

\[
\boxed{
A
=
\iint_D
\sqrt{EG-F^2}\,du\,dv.
}
\]

%%%%%%%%%%%%%%%%%%%%%%%%%%%%%%%%%%%%%%%%%%%%%%%%%%%%%%%%%%%%
\subsection{Worked Example: Metric on a Sphere}
\label{subsec:worked-sphere-metric}
%%%%%%%%%%%%%%%%%%%%%%%%%%%%%%%%%%%%%%%%%%%%%%%%%%%%%%%%%%%%

\begin{workedexamplebox}{First Fundamental Form of a Sphere}

For

\[
\mathbf{X}(\theta,\phi)
=
\left(
R\sin\phi\cos\theta,
R\sin\phi\sin\theta,
R\cos\phi
\right),
\]

we obtain

\[
E
=
\mathbf{X}_\theta\cdot\mathbf{X}_\theta
=
R^2\sin^2\phi,
\]

\[
F
=
\mathbf{X}_\theta\cdot\mathbf{X}_\phi
=
0,
\]

and

\[
G
=
\mathbf{X}_\phi\cdot\mathbf{X}_\phi
=
R^2.
\]

Therefore

\[
\boxed{
ds^2
=
R^2\sin^2\phi\,d\theta^2
+
R^2\,d\phi^2.
}
\]

The area element is

\[
\boxed{
dA
=
R^2\sin\phi\,d\theta\,d\phi.
}
\]
\end{workedexamplebox}

%%%%%%%%%%%%%%%%%%%%%%%%%%%%%%%%%%%%%%%%%%%%%%%%%%%%%%%%%%%%
\subsection{The Second Fundamental Form}
\label{subsec:second-fundamental-form}
%%%%%%%%%%%%%%%%%%%%%%%%%%%%%%%%%%%%%%%%%%%%%%%%%%%%%%%%%%%%

The first fundamental form describes intrinsic measurement.

The second fundamental form describes how the surface bends relative to the
ambient space.

Define

\[
e
=
\mathbf{X}_{uu}\cdot\mathbf{n},
\]

\[
f
=
\mathbf{X}_{uv}\cdot\mathbf{n},
\]

and

\[
g
=
\mathbf{X}_{vv}\cdot\mathbf{n}.
\]

Then

\[
\boxed{
II
=
e\,du^2
+
2f\,du\,dv
+
g\,dv^2.
}
\]

This is the \emph{second fundamental form}.

%%%%%%%%%%%%%%%%%%%%%%%%%%%%%%%%%%%%%%%%%%%%%%%%%%%%%%%%%%%%
\subsection{Normal Curvature}
\label{subsec:normal-curvature}
%%%%%%%%%%%%%%%%%%%%%%%%%%%%%%%%%%%%%%%%%%%%%%%%%%%%%%%%%%%%

At a point on a surface, different tangent directions may bend by different
amounts.

The curvature obtained by intersecting the surface with the plane generated
by a tangent direction and the surface normal is called the
\emph{normal curvature}.

If a tangent direction has coordinate differential \((du,dv)\), then

\[
\boxed{
\kappa_n
=
\frac{II}{I}
=
\frac{
e\,du^2+2f\,du\,dv+g\,dv^2
}{
E\,du^2+2F\,du\,dv+G\,dv^2
}.
}
\]

%%%%%%%%%%%%%%%%%%%%%%%%%%%%%%%%%%%%%%%%%%%%%%%%%%%%%%%%%%%%
\subsection{Principal Curvatures}
\label{subsec:principal-curvatures}
%%%%%%%%%%%%%%%%%%%%%%%%%%%%%%%%%%%%%%%%%%%%%%%%%%%%%%%%%%%%

Among all unit tangent directions, normal curvature attains extremal values.

\begin{definitionbox}{Principal Curvatures}
The maximum and minimum normal curvatures at a point are called the
\emph{principal curvatures} and are denoted

\[
\boxed{
\kappa_1,\qquad\kappa_2.
}
\]
\end{definitionbox}

The associated tangent directions are called \emph{principal directions}.

%%%%%%%%%%%%%%%%%%%%%%%%%%%%%%%%%%%%%%%%%%%%%%%%%%%%%%%%%%%%
\subsection{The Shape Operator}
\label{subsec:shape-operator}
%%%%%%%%%%%%%%%%%%%%%%%%%%%%%%%%%%%%%%%%%%%%%%%%%%%%%%%%%%%%

The unit normal

\[
\mathbf{n}
\]

changes as one moves along the surface.

Its differential determines a linear map on the tangent plane.

One common sign convention defines the \emph{shape operator} by

\[
\boxed{
S_p(v)=-D_v\mathbf{n}.
}
\]

The eigenvalues of \(S_p\) are the principal curvatures

\[
\boxed{
\kappa_1,\qquad\kappa_2.
}
\]

Different texts may adopt the opposite sign convention.

%%%%%%%%%%%%%%%%%%%%%%%%%%%%%%%%%%%%%%%%%%%%%%%%%%%%%%%%%%%%
\subsection{Gaussian Curvature}
\label{subsec:gaussian-curvature}
%%%%%%%%%%%%%%%%%%%%%%%%%%%%%%%%%%%%%%%%%%%%%%%%%%%%%%%%%%%%

\begin{definitionbox}{Gaussian Curvature}
The \emph{Gaussian curvature} of a surface is

\[
\boxed{
K=\kappa_1\kappa_2.
}
\]
\end{definitionbox}

Equivalently,

\[
\boxed{
K
=
\frac{eg-f^2}
{EG-F^2}.
}
\]

Gaussian curvature measures an important intrinsic aspect of surface
curvature.

%%%%%%%%%%%%%%%%%%%%%%%%%%%%%%%%%%%%%%%%%%%%%%%%%%%%%%%%%%%%
\subsection{Mean Curvature}
\label{subsec:mean-curvature}
%%%%%%%%%%%%%%%%%%%%%%%%%%%%%%%%%%%%%%%%%%%%%%%%%%%%%%%%%%%%

\begin{definitionbox}{Mean Curvature}
The \emph{mean curvature} is

\[
\boxed{
H
=
\frac{\kappa_1+\kappa_2}{2}.
}
\]
\end{definitionbox}

The letter

\[
H
\]

is read ``H.''

The sign of \(H\) depends on the choice of surface normal under the usual
signed-curvature convention.

Gaussian curvature does not change when the normal is reversed.

%%%%%%%%%%%%%%%%%%%%%%%%%%%%%%%%%%%%%%%%%%%%%%%%%%%%%%%%%%%%
\subsection{Classification by Gaussian Curvature}
\label{subsec:surface-point-classification}
%%%%%%%%%%%%%%%%%%%%%%%%%%%%%%%%%%%%%%%%%%%%%%%%%%%%%%%%%%%%

At a point:

\[
\boxed{
\begin{array}{c|c}
K>0 & \text{elliptic point}\\
K<0 & \text{hyperbolic point}\\
K=0 & \text{parabolic or planar behavior}
\end{array}
}
\]

Typical examples are:

\begin{itemize}
    \item sphere: \(K>0\);
    \item saddle surface: \(K<0\);
    \item cylinder: \(K=0\);
    \item plane: \(K=0\).
\end{itemize}

The case \(K=0\) requires additional information to distinguish planar and
parabolic points.

%%%%%%%%%%%%%%%%%%%%%%%%%%%%%%%%%%%%%%%%%%%%%%%%%%%%%%%%%%%%
\subsection{Curvature of a Sphere}
\label{subsec:curvature-sphere}
%%%%%%%%%%%%%%%%%%%%%%%%%%%%%%%%%%%%%%%%%%%%%%%%%%%%%%%%%%%%

For a sphere of radius \(R\), the two principal curvatures have magnitude

\[
\frac{1}{R}.
\]

Their signs depend on the orientation convention.

Therefore

\[
\boxed{
K=\frac{1}{R^2}.
}
\]

The mean curvature has magnitude

\[
\boxed{
|H|=\frac{1}{R}.
}
\]

Thus smaller spheres have greater curvature.

%%%%%%%%%%%%%%%%%%%%%%%%%%%%%%%%%%%%%%%%%%%%%%%%%%%%%%%%%%%%
\subsection{Curvature of a Cylinder}
\label{subsec:curvature-cylinder}
%%%%%%%%%%%%%%%%%%%%%%%%%%%%%%%%%%%%%%%%%%%%%%%%%%%%%%%%%%%%

For a circular cylinder of radius \(R\), one principal direction follows the
circular cross-section while the other follows the axis.

The principal curvatures have magnitudes

\[
\frac{1}{R}
\]

and

\[
0.
\]

Therefore

\[
\boxed{
K=0.
}
\]

This is surprising at first because a cylinder visibly bends in
three-dimensional space.

%%%%%%%%%%%%%%%%%%%%%%%%%%%%%%%%%%%%%%%%%%%%%%%%%%%%%%%%%%%%
\subsection{Intrinsic and Extrinsic Curvature}
\label{subsec:intrinsic-extrinsic-curvature}
%%%%%%%%%%%%%%%%%%%%%%%%%%%%%%%%%%%%%%%%%%%%%%%%%%%%%%%%%%%%

The cylinder demonstrates the distinction between two ideas.

\begin{definitionbox}{Extrinsic Geometry}
\emph{Extrinsic geometry} describes how a geometric object sits or bends
inside an ambient space.
\end{definitionbox}

\begin{definitionbox}{Intrinsic Geometry}
\emph{Intrinsic geometry} consists of geometric information measurable
entirely within the object itself.
\end{definitionbox}

A sheet of paper can be rolled into a cylinder without stretching it.

The extrinsic appearance changes, but its intrinsic metric remains locally
Euclidean.

%%%%%%%%%%%%%%%%%%%%%%%%%%%%%%%%%%%%%%%%%%%%%%%%%%%%%%%%%%%%
\subsection{Gauss's Theorema Egregium}
\label{subsec:theorema-egregium}
%%%%%%%%%%%%%%%%%%%%%%%%%%%%%%%%%%%%%%%%%%%%%%%%%%%%%%%%%%%%

One of the central discoveries of differential geometry is that Gaussian
curvature is intrinsic.

\begin{theorem}[Gauss's Theorema Egregium]
The Gaussian curvature of a surface is determined entirely by its first
fundamental form and its derivatives.

Consequently, Gaussian curvature is preserved by local isometries.
\end{theorem}

\emph{Theorema Egregium} means ``remarkable theorem.''

Although

\[
K=\kappa_1\kappa_2
\]

was initially defined using the surface normal and ambient space, the final
quantity can be determined intrinsically.

%%%%%%%%%%%%%%%%%%%%%%%%%%%%%%%%%%%%%%%%%%%%%%%%%%%%%%%%%%%%
\subsection{Why a Flat Map Distorts Earth}
\label{subsec:flat-map-distortion}
%%%%%%%%%%%%%%%%%%%%%%%%%%%%%%%%%%%%%%%%%%%%%%%%%%%%%%%%%%%%

A sphere has

\[
K=\frac{1}{R^2}>0,
\]

while the Euclidean plane has

\[
K=0.
\]

Because Gaussian curvature is preserved by local isometries, no map from a
spherical region to the plane can preserve all distances locally over an
extended open region.

Therefore every flat map of Earth must introduce some form of distortion.

Depending on the projection, one may sacrifice:

\begin{itemize}
    \item distance;
    \item area;
    \item angles;
    \item shapes;
    \item directions.
\end{itemize}

%%%%%%%%%%%%%%%%%%%%%%%%%%%%%%%%%%%%%%%%%%%%%%%%%%%%%%%%%%%%
\subsection{Minimal Surfaces}
\label{subsec:minimal-surfaces}
%%%%%%%%%%%%%%%%%%%%%%%%%%%%%%%%%%%%%%%%%%%%%%%%%%%%%%%%%%%%

\begin{definitionbox}{Minimal Surface}
A surface is \emph{minimal} when its mean curvature satisfies

\[
\boxed{H=0}
\]

at every point.
\end{definitionbox}

Examples include:

\begin{itemize}
    \item planes;
    \item catenoids;
    \item helicoids.
\end{itemize}

Minimal surfaces arise naturally as critical points of the area functional.

Soap films provide important physical models.

%%%%%%%%%%%%%%%%%%%%%%%%%%%%%%%%%%%%%%%%%%%%%%%%%%%%%%%%%%%%
\subsection{The Catenoid}
\label{subsec:catenoid}
%%%%%%%%%%%%%%%%%%%%%%%%%%%%%%%%%%%%%%%%%%%%%%%%%%%%%%%%%%%%

A standard parametrization of a catenoid is

\[
\boxed{
\mathbf{X}(u,v)
=
\left(
a\cosh v\cos u,
a\cosh v\sin u,
av
\right).
}
\]

The catenoid is a surface of revolution.

Its mean curvature satisfies

\[
\boxed{H=0.}
\]

It is one of the classical nonplanar minimal surfaces.

%%%%%%%%%%%%%%%%%%%%%%%%%%%%%%%%%%%%%%%%%%%%%%%%%%%%%%%%%%%%
\subsection{Geodesics on Surfaces}
\label{subsec:surface-geodesics}
%%%%%%%%%%%%%%%%%%%%%%%%%%%%%%%%%%%%%%%%%%%%%%%%%%%%%%%%%%%%

A geodesic generalizes the notion of a straight line to a curved surface or
manifold.

Informally, a geodesic follows the straightest possible path allowed by the
intrinsic geometry.

For a surface embedded in \(\R^3\), a unit-speed curve is geodesic when its
acceleration has no tangential component.

Thus

\[
\boxed{
\gamma''(s)
\perp
T_{\gamma(s)}S.
}
\]

Its acceleration points entirely in the normal direction.

%%%%%%%%%%%%%%%%%%%%%%%%%%%%%%%%%%%%%%%%%%%%%%%%%%%%%%%%%%%%
\subsection{Examples of Geodesics}
\label{subsec:geodesic-examples}
%%%%%%%%%%%%%%%%%%%%%%%%%%%%%%%%%%%%%%%%%%%%%%%%%%%%%%%%%%%%

Examples include:

\[
\begin{array}{c|c}
\textbf{Space} & \textbf{Geodesics}\\
\hline
\R^2 & \text{straight lines}\\
S^2 & \text{great circles}\\
\mathbb{H}^2 & \text{hyperbolic lines}\\
\text{cylinder} & \text{locally straight unrolled paths}
\end{array}
\]

On a cylinder, a helical geodesic becomes a straight line when the cylinder
is unrolled onto a plane.

%%%%%%%%%%%%%%%%%%%%%%%%%%%%%%%%%%%%%%%%%%%%%%%%%%%%%%%%%%%%
\subsection{Geodesics and Shortest Paths}
\label{subsec:geodesics-shortest-paths}
%%%%%%%%%%%%%%%%%%%%%%%%%%%%%%%%%%%%%%%%%%%%%%%%%%%%%%%%%%%%

Geodesics are locally length minimizing under suitable conditions, but a
geodesic need not remain globally shortest.

For example, on a sphere, a great-circle arc longer than half a great circle
is still part of a geodesic, but the opposite arc provides a shorter route
between the same endpoints.

Thus:

\[
\boxed{
\text{geodesic}
\not\Longleftrightarrow
\text{globally shortest path}.
}
\]

The distinction between local and global minimization is important.

%%%%%%%%%%%%%%%%%%%%%%%%%%%%%%%%%%%%%%%%%%%%%%%%%%%%%%%%%%%%
\subsection{The Energy Functional}
\label{subsec:geodesic-energy}
%%%%%%%%%%%%%%%%%%%%%%%%%%%%%%%%%%%%%%%%%%%%%%%%%%%%%%%%%%%%

For a curve

\[
\gamma:[a,b]\to M
\]

on a Riemannian manifold, its energy is

\[
\boxed{
E(\gamma)
=
\frac12
\int_a^b
\|\gamma'(t)\|^2\,dt.
}
\]

Geodesics arise as critical points of this functional.

For curves with fixed endpoints and suitable parametrization, minimizing
energy is closely related to minimizing length.

This creates a bridge between differential geometry and the Calculus of
Variations.

%%%%%%%%%%%%%%%%%%%%%%%%%%%%%%%%%%%%%%%%%%%%%%%%%%%%%%%%%%%%
\subsection{From Surfaces to Manifolds}
\label{subsec:surfaces-to-manifolds}
%%%%%%%%%%%%%%%%%%%%%%%%%%%%%%%%%%%%%%%%%%%%%%%%%%%%%%%%%%%%

A surface is locally two-dimensional.

Differential geometry generalizes this idea to arbitrary dimension.

\begin{definitionbox}{Smooth Manifold}
An \(n\)-dimensional \emph{smooth manifold} is a space that locally resembles
\(\R^n\) and whose overlapping coordinate systems are related by smooth
transition maps.
\end{definitionbox}

Examples include:

\[
\R^n,
\qquad
S^n,
\qquad
T^n,
\qquad
\mathbb{P}^n(\R),
\]

under their standard smooth structures.

%%%%%%%%%%%%%%%%%%%%%%%%%%%%%%%%%%%%%%%%%%%%%%%%%%%%%%%%%%%%
\subsection{Coordinate Charts}
\label{subsec:differential-coordinate-charts}
%%%%%%%%%%%%%%%%%%%%%%%%%%%%%%%%%%%%%%%%%%%%%%%%%%%%%%%%%%%%

A coordinate chart assigns ordinary Euclidean coordinates to a portion of a
manifold.

A chart is typically written

\[
\boxed{
\varphi:U\to\varphi(U)\subseteq\R^n.
}
\]

The Greek letter

\[
\varphi
\]

is read ``phi'' or ``varphi.''

The set \(U\) is an open region of the manifold.

The coordinate map allows multivariable calculus to be performed locally.

%%%%%%%%%%%%%%%%%%%%%%%%%%%%%%%%%%%%%%%%%%%%%%%%%%%%%%%%%%%%
\subsection{Atlases}
\label{subsec:manifold-atlas}
%%%%%%%%%%%%%%%%%%%%%%%%%%%%%%%%%%%%%%%%%%%%%%%%%%%%%%%%%%%%

A collection of compatible coordinate charts covering the manifold is
called an \emph{atlas}.

If

\[
(U,\varphi)
\]

and

\[
(V,\psi)
\]

overlap, then the transition map is

\[
\boxed{
\psi\circ\varphi^{-1}.
}
\]

For a smooth atlas, these transition maps must be smooth wherever defined.

%%%%%%%%%%%%%%%%%%%%%%%%%%%%%%%%%%%%%%%%%%%%%%%%%%%%%%%%%%%%
\subsection{Why Multiple Charts Are Necessary}
\label{subsec:why-multiple-charts}
%%%%%%%%%%%%%%%%%%%%%%%%%%%%%%%%%%%%%%%%%%%%%%%%%%%%%%%%%%%%

A sphere cannot be represented by one ordinary smooth coordinate system
without singularities or omissions.

Longitude and latitude, for example, become problematic at the poles.

Instead, several overlapping coordinate charts are used.

This illustrates a recurring principle:

\[
\boxed{
\text{manifolds are locally Euclidean but need not be globally Euclidean.}
}
\]

%%%%%%%%%%%%%%%%%%%%%%%%%%%%%%%%%%%%%%%%%%%%%%%%%%%%%%%%%%%%
\subsection{Tangent Spaces}
\label{subsec:manifold-tangent-spaces}
%%%%%%%%%%%%%%%%%%%%%%%%%%%%%%%%%%%%%%%%%%%%%%%%%%%%%%%%%%%%

At each point

\[
p\in M,
\]

an \(n\)-dimensional smooth manifold has an \(n\)-dimensional tangent vector
space

\[
\boxed{
T_pM.
}
\]

For a surface in \(\R^3\), this agrees with the familiar tangent plane.

For an abstract manifold, the tangent space is defined intrinsically without
requiring an ambient Euclidean space.

%%%%%%%%%%%%%%%%%%%%%%%%%%%%%%%%%%%%%%%%%%%%%%%%%%%%%%%%%%%%
\subsection{Tangent Vectors from Curves}
\label{subsec:tangent-vectors-curves}
%%%%%%%%%%%%%%%%%%%%%%%%%%%%%%%%%%%%%%%%%%%%%%%%%%%%%%%%%%%%

A tangent vector at \(p\) may be represented by the velocity of a smooth
curve through \(p\).

If

\[
\gamma(0)=p,
\]

then

\[
\boxed{
\gamma'(0)\in T_pM.
}
\]

Different curves may determine the same tangent vector.

This gives a geometric interpretation of tangent vectors independent of a
particular embedding.

%%%%%%%%%%%%%%%%%%%%%%%%%%%%%%%%%%%%%%%%%%%%%%%%%%%%%%%%%%%%
\subsection{Coordinate Basis}
\label{subsec:coordinate-basis}
%%%%%%%%%%%%%%%%%%%%%%%%%%%%%%%%%%%%%%%%%%%%%%%%%%%%%%%%%%%%

In local coordinates

\[
(x^1,\ldots,x^n),
\]

a standard basis for the tangent space is written

\[
\boxed{
\frac{\partial}{\partial x^1},
\ldots,
\frac{\partial}{\partial x^n}.
}
\]

A tangent vector may therefore be written

\[
\boxed{
v
=
v^i
\frac{\partial}{\partial x^i}.
}
\]

Here repeated-index summation is often understood.

This convention is called the \emph{Einstein summation convention}.

%%%%%%%%%%%%%%%%%%%%%%%%%%%%%%%%%%%%%%%%%%%%%%%%%%%%%%%%%%%%
\subsection{Einstein Summation Convention}
\label{subsec:einstein-summation}
%%%%%%%%%%%%%%%%%%%%%%%%%%%%%%%%%%%%%%%%%%%%%%%%%%%%%%%%%%%%

Under the Einstein convention,

\[
v^i
\frac{\partial}{\partial x^i}
\]

means

\[
\sum_{i=1}^{n}
v^i
\frac{\partial}{\partial x^i}.
\]

Likewise,

\[
g_{ij}v^iw^j
\]

means

\[
\sum_{i=1}^{n}
\sum_{j=1}^{n}
g_{ij}v^iw^j.
\]

This convention greatly reduces notation in tensor calculations.

%%%%%%%%%%%%%%%%%%%%%%%%%%%%%%%%%%%%%%%%%%%%%%%%%%%%%%%%%%%%
\subsection{Cotangent Spaces}
\label{subsec:cotangent-spaces}
%%%%%%%%%%%%%%%%%%%%%%%%%%%%%%%%%%%%%%%%%%%%%%%%%%%%%%%%%%%%

The dual vector space of

\[
T_pM
\]

is the \emph{cotangent space}

\[
\boxed{
T_p^*M.
}
\]

Its elements are linear functionals on tangent vectors.

The coordinate basis is

\[
\boxed{
dx^1,\ldots,dx^n.
}
\]

They satisfy

\[
\boxed{
dx^i
\left(
\frac{\partial}{\partial x^j}
\right)
=
\delta^i_j.
}
\]

The symbol

\[
\delta^i_j
\]

is the Kronecker delta.

%%%%%%%%%%%%%%%%%%%%%%%%%%%%%%%%%%%%%%%%%%%%%%%%%%%%%%%%%%%%
\subsection{Riemannian Metrics}
\label{subsec:riemannian-metrics}
%%%%%%%%%%%%%%%%%%%%%%%%%%%%%%%%%%%%%%%%%%%%%%%%%%%%%%%%%%%%

To measure lengths and angles on a manifold, each tangent space needs an
inner product.

\begin{definitionbox}{Riemannian Metric}
A \emph{Riemannian metric} assigns to every point \(p\in M\) an inner
product

\[
\boxed{
g_p:T_pM\times T_pM\to\R
}
\]

that varies smoothly with \(p\).
\end{definitionbox}

In coordinates,

\[
\boxed{
g
=
g_{ij}\,dx^i\,dx^j.
}
\]

The matrix

\[
(g_{ij})
\]

is symmetric and positive definite.

%%%%%%%%%%%%%%%%%%%%%%%%%%%%%%%%%%%%%%%%%%%%%%%%%%%%%%%%%%%%
\subsection{Length from a Riemannian Metric}
\label{subsec:riemannian-length}
%%%%%%%%%%%%%%%%%%%%%%%%%%%%%%%%%%%%%%%%%%%%%%%%%%%%%%%%%%%%

If

\[
\gamma(t)
\]

is a smooth curve on \(M\), its speed is

\[
\boxed{
\|\gamma'(t)\|_g
=
\sqrt{
g_{\gamma(t)}
\bigl(
\gamma'(t),
\gamma'(t)
\bigr)
}.
}
\]

Its length is

\[
\boxed{
L(\gamma)
=
\int_a^b
\|\gamma'(t)\|_g\,dt.
}
\]

Thus the metric tensor generalizes the Euclidean dot product.

%%%%%%%%%%%%%%%%%%%%%%%%%%%%%%%%%%%%%%%%%%%%%%%%%%%%%%%%%%%%
\subsection{Geodesic Distance}
\label{subsec:geodesic-distance}
%%%%%%%%%%%%%%%%%%%%%%%%%%%%%%%%%%%%%%%%%%%%%%%%%%%%%%%%%%%%

The Riemannian distance between two points is defined by minimizing curve
length:

\[
\boxed{
d(p,q)
=
\inf_{\gamma}
L(\gamma),
}
\]

where the infimum is taken over suitable smooth or piecewise smooth curves
joining \(p\) to \(q\).

The symbol

\[
\inf
\]

means ``infimum,'' or greatest lower bound.

Thus the metric tensor generates a metric-space distance under standard
connectedness assumptions.

%%%%%%%%%%%%%%%%%%%%%%%%%%%%%%%%%%%%%%%%%%%%%%%%%%%%%%%%%%%%
\subsection{Connections and Covariant Derivatives}
\label{subsec:connections-covariant-derivatives}
%%%%%%%%%%%%%%%%%%%%%%%%%%%%%%%%%%%%%%%%%%%%%%%%%%%%%%%%%%%%

Vectors at different points lie in different tangent spaces.

Therefore ordinary subtraction of tangent vectors at different points is not
intrinsically meaningful.

A \emph{connection} provides a method for differentiating vector fields while
respecting the manifold structure.

The corresponding derivative is written

\[
\boxed{
\nabla_XY.
}
\]

The symbol

\[
\nabla
\]

is called ``nabla'' or ``del.''

Here

\[
\nabla_XY
\]

is read ``the covariant derivative of \(Y\) in the direction \(X\).''

%%%%%%%%%%%%%%%%%%%%%%%%%%%%%%%%%%%%%%%%%%%%%%%%%%%%%%%%%%%%
\subsection{The Levi-Civita Connection}
\label{subsec:levi-civita-connection}
%%%%%%%%%%%%%%%%%%%%%%%%%%%%%%%%%%%%%%%%%%%%%%%%%%%%%%%%%%%%

A Riemannian metric determines a distinguished connection.

\begin{theorem}[Fundamental Theorem of Riemannian Geometry]
Every Riemannian manifold has a unique connection that is

\begin{enumerate}
    \item compatible with the metric; and
    \item torsion-free.
\end{enumerate}

This connection is called the \emph{Levi-Civita connection}.
\end{theorem}

This theorem is one of the foundations of Riemannian geometry.

%%%%%%%%%%%%%%%%%%%%%%%%%%%%%%%%%%%%%%%%%%%%%%%%%%%%%%%%%%%%
\subsection{Christoffel Symbols}
\label{subsec:christoffel-symbols}
%%%%%%%%%%%%%%%%%%%%%%%%%%%%%%%%%%%%%%%%%%%%%%%%%%%%%%%%%%%%

In local coordinates, the Levi-Civita connection is described by
\emph{Christoffel symbols}

\[
\boxed{
\Gamma^k_{ij}.
}
\]

The Greek capital letter

\[
\Gamma
\]

is read ``capital gamma.''

They are given by

\[
\boxed{
\Gamma^k_{ij}
=
\frac12
g^{k\ell}
\left(
\frac{\partial g_{j\ell}}{\partial x^i}
+
\frac{\partial g_{i\ell}}{\partial x^j}
-
\frac{\partial g_{ij}}{\partial x^\ell}
\right).
}
\]

Here

\[
(g^{ij})
\]

is the inverse matrix of

\[
(g_{ij}).
\]

%%%%%%%%%%%%%%%%%%%%%%%%%%%%%%%%%%%%%%%%%%%%%%%%%%%%%%%%%%%%
\subsection{The Geodesic Equation}
\label{subsec:geodesic-equation}
%%%%%%%%%%%%%%%%%%%%%%%%%%%%%%%%%%%%%%%%%%%%%%%%%%%%%%%%%%%%

A curve

\[
x^k=x^k(t)
\]

is a geodesic when its tangent vector is parallel transported along itself:

\[
\boxed{
\nabla_{\dot{\gamma}}\dot{\gamma}=0.
}
\]

In coordinates, this becomes

\[
\boxed{
\frac{d^2x^k}{dt^2}
+
\Gamma^k_{ij}
\frac{dx^i}{dt}
\frac{dx^j}{dt}
=
0.
}
\]

This is a nonlinear system of second-order differential equations.

Thus geodesics connect geometry directly with differential equations.

%%%%%%%%%%%%%%%%%%%%%%%%%%%%%%%%%%%%%%%%%%%%%%%%%%%%%%%%%%%%
\subsection{Worked Example: Euclidean Geodesics}
\label{subsec:worked-euclidean-geodesics}
%%%%%%%%%%%%%%%%%%%%%%%%%%%%%%%%%%%%%%%%%%%%%%%%%%%%%%%%%%%%

\begin{workedexamplebox}{Geodesics in Cartesian Coordinates}

In Euclidean Cartesian coordinates,

\[
g_{ij}=\delta_{ij}.
\]

The metric coefficients are constant, so

\[
\frac{\partial g_{ij}}{\partial x^k}=0.
\]

Therefore

\[
\Gamma^k_{ij}=0.
\]

The geodesic equation becomes

\[
\frac{d^2x^k}{dt^2}=0.
\]

Integrating twice,

\[
x^k(t)=a^kt+b^k.
\]

Thus Euclidean geodesics are straight lines.

\begin{answerbox}
\[
\boxed{
\gamma(t)=\mathbf{a}t+\mathbf{b}
}
\]
\end{answerbox}
\end{workedexamplebox}

%%%%%%%%%%%%%%%%%%%%%%%%%%%%%%%%%%%%%%%%%%%%%%%%%%%%%%%%%%%%
\subsection{Parallel Transport}
\label{subsec:parallel-transport}
%%%%%%%%%%%%%%%%%%%%%%%%%%%%%%%%%%%%%%%%%%%%%%%%%%%%%%%%%%%%

A vector field \(V(t)\) along a curve \(\gamma(t)\) is parallel when

\[
\boxed{
\nabla_{\dot{\gamma}}V=0.
}
\]

This process is called \emph{parallel transport}.

In Euclidean space, parallel transport preserves an ordinary fixed vector
direction.

On a curved manifold, transporting a vector around a closed loop may return
a vector with a different direction.

%%%%%%%%%%%%%%%%%%%%%%%%%%%%%%%%%%%%%%%%%%%%%%%%%%%%%%%%%%%%
\subsection{Curvature from Parallel Transport}
\label{subsec:curvature-parallel-transport}
%%%%%%%%%%%%%%%%%%%%%%%%%%%%%%%%%%%%%%%%%%%%%%%%%%%%%%%%%%%%

The failure of parallel transport to return vectors unchanged around small
loops provides an intrinsic manifestation of curvature.

Schematically,

\[
\boxed{
\text{curvature}
\quad\Longleftrightarrow\quad
\text{path dependence of parallel transport}.
}
\]

On a flat Euclidean plane, parallel transport around a contractible closed
loop returns the vector unchanged.

On a sphere, the vector may rotate.

%%%%%%%%%%%%%%%%%%%%%%%%%%%%%%%%%%%%%%%%%%%%%%%%%%%%%%%%%%%%
\subsection{The Riemann Curvature Tensor}
\label{subsec:riemann-curvature-tensor}
%%%%%%%%%%%%%%%%%%%%%%%%%%%%%%%%%%%%%%%%%%%%%%%%%%%%%%%%%%%%

The curvature of a Riemannian manifold is encoded by the
\emph{Riemann curvature tensor}.

One common sign convention defines

\[
\boxed{
R(X,Y)Z
=
\nabla_X\nabla_YZ
-
\nabla_Y\nabla_XZ
-
\nabla_{[X,Y]}Z.
}
\]

The notation

\[
[X,Y]
\]

denotes the Lie bracket of vector fields.

Different references may use the opposite overall sign convention for
\(R\).

%%%%%%%%%%%%%%%%%%%%%%%%%%%%%%%%%%%%%%%%%%%%%%%%%%%%%%%%%%%%
\subsection{Sectional Curvature}
\label{subsec:sectional-curvature}
%%%%%%%%%%%%%%%%%%%%%%%%%%%%%%%%%%%%%%%%%%%%%%%%%%%%%%%%%%%%

In dimensions greater than two, curvature may depend on the chosen
two-dimensional tangent direction.

Given a two-dimensional plane

\[
\sigma\subseteq T_pM,
\]

its \emph{sectional curvature} is denoted

\[
\boxed{K(\sigma).}
\]

For a two-dimensional Riemannian manifold, there is only one tangent
two-plane at each point, and sectional curvature reduces to Gaussian
curvature.

%%%%%%%%%%%%%%%%%%%%%%%%%%%%%%%%%%%%%%%%%%%%%%%%%%%%%%%%%%%%
\subsection{Constant Curvature Spaces}
\label{subsec:constant-curvature-spaces}
%%%%%%%%%%%%%%%%%%%%%%%%%%%%%%%%%%%%%%%%%%%%%%%%%%%%%%%%%%%%

The standard model spaces satisfy:

\[
\boxed{
\begin{array}{c|c}
\textbf{Space} & \textbf{Sectional curvature}\\
\hline
S^n_R & +\frac{1}{R^2}\\[2mm]
\R^n & 0\\[2mm]
\mathbb{H}^n_R & -\frac{1}{R^2}
\end{array}
}
\]

These generalize the spherical, Euclidean, and hyperbolic geometries studied
in \cref{sec:non-euclidean-geometry}.

%%%%%%%%%%%%%%%%%%%%%%%%%%%%%%%%%%%%%%%%%%%%%%%%%%%%%%%%%%%%
\subsection{Ricci Curvature}
\label{subsec:ricci-curvature}
%%%%%%%%%%%%%%%%%%%%%%%%%%%%%%%%%%%%%%%%%%%%%%%%%%%%%%%%%%%%

The Riemann tensor contains a large amount of information.

A contraction of this tensor produces the \emph{Ricci curvature}.

It is represented by the Ricci tensor

\[
\boxed{
\operatorname{Ric}_{ij}.
}
\]

Conceptually, Ricci curvature measures how volumes and families of geodesics
behave relative to Euclidean expectations.

It plays a central role in geometry and general relativity.

%%%%%%%%%%%%%%%%%%%%%%%%%%%%%%%%%%%%%%%%%%%%%%%%%%%%%%%%%%%%
\subsection{Scalar Curvature}
\label{subsec:scalar-curvature}
%%%%%%%%%%%%%%%%%%%%%%%%%%%%%%%%%%%%%%%%%%%%%%%%%%%%%%%%%%%%

Contracting the Ricci tensor again produces the \emph{scalar curvature}

\[
\boxed{
R
=
g^{ij}\operatorname{Ric}_{ij}.
}
\]

The same letter \(R\) is often used for both the Riemann tensor and scalar
curvature in different contexts, so the intended meaning must be determined
from notation and context.

Scalar curvature compresses curvature information into a single number at
each point.

%%%%%%%%%%%%%%%%%%%%%%%%%%%%%%%%%%%%%%%%%%%%%%%%%%%%%%%%%%%%
\subsection{Curvature Hierarchy}
\label{subsec:curvature-hierarchy}
%%%%%%%%%%%%%%%%%%%%%%%%%%%%%%%%%%%%%%%%%%%%%%%%%%%%%%%%%%%%

A useful conceptual hierarchy is

\begin{centereddiagram}
\begin{tikzpicture}[
node distance=9mm,
every node/.style={
draw,
rounded corners,
minimum width=5.2cm,
minimum height=8mm,
align=center
}
]
\node (r) {Riemann curvature tensor};
\node (ric) [below=of r] {Ricci curvature};
\node (sc) [below=of ric] {Scalar curvature};

\draw[-{Latex[length=3mm]},thick] (r)--node[right]{contraction}(ric);
\draw[-{Latex[length=3mm]},thick] (ric)--node[right]{contraction}(sc);
\end{tikzpicture}
\end{centereddiagram}

Each contraction loses information but produces quantities useful for
different geometric questions.

%%%%%%%%%%%%%%%%%%%%%%%%%%%%%%%%%%%%%%%%%%%%%%%%%%%%%%%%%%%%
\subsection{Exponential Map}
\label{subsec:exponential-map}
%%%%%%%%%%%%%%%%%%%%%%%%%%%%%%%%%%%%%%%%%%%%%%%%%%%%%%%%%%%%

Given

\[
p\in M
\]

and a tangent vector

\[
v\in T_pM,
\]

there is, locally, a geodesic satisfying

\[
\gamma(0)=p,
\qquad
\gamma'(0)=v.
\]

The \emph{exponential map} is defined by

\[
\boxed{
\exp_p(v)=\gamma_v(1),
}
\]

whenever the geodesic exists up to parameter \(1\).

Despite its name, this is not generally the ordinary exponential function.

%%%%%%%%%%%%%%%%%%%%%%%%%%%%%%%%%%%%%%%%%%%%%%%%%%%%%%%%%%%%
\subsection{Normal Coordinates}
\label{subsec:normal-coordinates}
%%%%%%%%%%%%%%%%%%%%%%%%%%%%%%%%%%%%%%%%%%%%%%%%%%%%%%%%%%%%

Near a point \(p\), the exponential map can be used to construct
\emph{normal coordinates}.

In suitable normal coordinates centered at \(p\),

\[
\boxed{
g_{ij}(p)=\delta_{ij}
}
\]

and

\[
\boxed{
\Gamma^k_{ij}(p)=0.
}
\]

Thus the geometry looks Euclidean to first order at a point.

Curvature appears in higher-order behavior.

%%%%%%%%%%%%%%%%%%%%%%%%%%%%%%%%%%%%%%%%%%%%%%%%%%%%%%%%%%%%
\subsection{Geodesic Completeness}
\label{subsec:geodesic-completeness}
%%%%%%%%%%%%%%%%%%%%%%%%%%%%%%%%%%%%%%%%%%%%%%%%%%%%%%%%%%%%

A Riemannian manifold is \emph{geodesically complete} if every geodesic can
be extended indefinitely in both parameter directions.

Examples include:

\[
\R^n,
\qquad
S^n,
\qquad
\mathbb{H}^n
\]

with their standard complete metrics.

Completeness links local differential equations with global geometry.

%%%%%%%%%%%%%%%%%%%%%%%%%%%%%%%%%%%%%%%%%%%%%%%%%%%%%%%%%%%%
\subsection{Hopf--Rinow Theorem}
\label{subsec:hopf-rinow}
%%%%%%%%%%%%%%%%%%%%%%%%%%%%%%%%%%%%%%%%%%%%%%%%%%%%%%%%%%%%

\begin{theorem}[Hopf--Rinow Theorem --- Conceptual Form]
For a connected Riemannian manifold, several natural notions of completeness
are equivalent.

In particular, metric completeness is equivalent to geodesic completeness,
and in the complete case any two points can be joined by a minimizing
geodesic.
\end{theorem}

This theorem is a fundamental local-to-global result in Riemannian geometry.

%%%%%%%%%%%%%%%%%%%%%%%%%%%%%%%%%%%%%%%%%%%%%%%%%%%%%%%%%%%%
\subsection{Jacobi Fields}
\label{subsec:jacobi-fields}
%%%%%%%%%%%%%%%%%%%%%%%%%%%%%%%%%%%%%%%%%%%%%%%%%%%%%%%%%%%%

To understand how nearby geodesics separate, one studies \emph{Jacobi
fields}.

A Jacobi field \(J\) along a geodesic \(\gamma\) satisfies an equation of the
form

\[
\boxed{
\frac{D^2J}{dt^2}
+
R(J,\dot{\gamma})\dot{\gamma}
=
0.
}
\]

Here

\[
\frac{D}{dt}
\]

denotes covariant differentiation along the curve.

The equation shows directly how curvature controls the behavior of nearby
geodesics.

%%%%%%%%%%%%%%%%%%%%%%%%%%%%%%%%%%%%%%%%%%%%%%%%%%%%%%%%%%%%
\subsection{Curvature and Geodesic Separation}
\label{subsec:curvature-geodesic-separation}
%%%%%%%%%%%%%%%%%%%%%%%%%%%%%%%%%%%%%%%%%%%%%%%%%%%%%%%%%%%%

Very roughly:

\[
\boxed{
\begin{array}{c|c}
K>0 & \text{nearby geodesics tend to reconverge}\\
K=0 & \text{nearby geodesics behave linearly}\\
K<0 & \text{nearby geodesics tend to separate}
\end{array}
}
\]

This helps explain why hyperbolic space expands much more rapidly than
Euclidean space.

The precise behavior depends on sectional curvature and the Jacobi equation.

%%%%%%%%%%%%%%%%%%%%%%%%%%%%%%%%%%%%%%%%%%%%%%%%%%%%%%%%%%%%
\subsection{Differential Forms}
\label{subsec:differential-forms-geometry}
%%%%%%%%%%%%%%%%%%%%%%%%%%%%%%%%%%%%%%%%%%%%%%%%%%%%%%%%%%%%

Differential forms provide a coordinate-independent language for integration
on manifolds.

A differential \(k\)-form is an alternating covariant tensor field.

Examples include:

\[
f
\]

as a \(0\)-form,

\[
\omega=P\,dx+Q\,dy
\]

as a \(1\)-form, and

\[
\omega=f\,dx\wedge dy
\]

as a \(2\)-form.

The symbol

\[
\wedge
\]

is called the \emph{wedge product}.

The systematic theory belongs to Differential Topology and Vector Calculus.

%%%%%%%%%%%%%%%%%%%%%%%%%%%%%%%%%%%%%%%%%%%%%%%%%%%%%%%%%%%%
\subsection{Exterior Derivative}
\label{subsec:exterior-derivative}
%%%%%%%%%%%%%%%%%%%%%%%%%%%%%%%%%%%%%%%%%%%%%%%%%%%%%%%%%%%%

The exterior derivative is an operator

\[
\boxed{
d:\Omega^k(M)\to\Omega^{k+1}(M).
}
\]

It satisfies

\[
\boxed{d^2=0.}
\]

For a function \(f\),

\[
df
\]

is its differential.

For a \(1\)-form in the plane,

\[
\omega=P\,dx+Q\,dy,
\]

we have

\[
\boxed{
d\omega
=
\left(
\frac{\partial Q}{\partial x}
-
\frac{\partial P}{\partial y}
\right)
dx\wedge dy.
}
\]

%%%%%%%%%%%%%%%%%%%%%%%%%%%%%%%%%%%%%%%%%%%%%%%%%%%%%%%%%%%%
\subsection{Stokes' Theorem on Manifolds}
\label{subsec:stokes-manifolds}
%%%%%%%%%%%%%%%%%%%%%%%%%%%%%%%%%%%%%%%%%%%%%%%%%%%%%%%%%%%%

The classical fundamental theorem of calculus, Green's theorem, divergence
theorem, and Stokes' theorem are unified by the generalized Stokes theorem.

\begin{theorem}[Generalized Stokes Theorem]
For a suitably oriented manifold \(M\) with boundary and an appropriate
differential form \(\omega\),

\[
\boxed{
\int_M d\omega
=
\int_{\partial M}\omega.
}
\]
\end{theorem}

The notation

\[
\partial M
\]

means ``the boundary of \(M\).''

This theorem is one of the central bridges between analysis, geometry, and
topology.

%%%%%%%%%%%%%%%%%%%%%%%%%%%%%%%%%%%%%%%%%%%%%%%%%%%%%%%%%%%%
\subsection{Gauss--Bonnet Theorem}
\label{subsec:gauss-bonnet}
%%%%%%%%%%%%%%%%%%%%%%%%%%%%%%%%%%%%%%%%%%%%%%%%%%%%%%%%%%%%

For a compact oriented surface without boundary, the Gauss--Bonnet theorem
states

\[
\boxed{
\int_M K\,dA
=
2\pi\chi(M).
}
\]

Here

\[
K
\]

is Gaussian curvature and

\[
\chi(M)
\]

is the Euler characteristic.

This remarkable formula relates a geometric quantity,

\[
\int_MK\,dA,
\]

to a topological invariant,

\[
\chi(M).
\]

%%%%%%%%%%%%%%%%%%%%%%%%%%%%%%%%%%%%%%%%%%%%%%%%%%%%%%%%%%%%
\subsection{Worked Example: Gauss--Bonnet on a Sphere}
\label{subsec:worked-gauss-bonnet-sphere}
%%%%%%%%%%%%%%%%%%%%%%%%%%%%%%%%%%%%%%%%%%%%%%%%%%%%%%%%%%%%

\begin{workedexamplebox}{Total Curvature of a Sphere}

For a sphere of radius \(R\),

\[
K=\frac{1}{R^2}.
\]

Its area is

\[
A=4\pi R^2.
\]

Therefore

\[
\int_{S^2}K\,dA
=
\frac{1}{R^2}
(4\pi R^2)
=
4\pi.
\]

Since

\[
\chi(S^2)=2,
\]

Gauss--Bonnet gives

\[
2\pi\chi(S^2)
=
2\pi(2)
=
4\pi.
\]

\begin{answerbox}
\[
\boxed{
\int_{S^2}K\,dA=4\pi
}
\]
\end{answerbox}
\end{workedexamplebox}

%%%%%%%%%%%%%%%%%%%%%%%%%%%%%%%%%%%%%%%%%%%%%%%%%%%%%%%%%%%%
\subsection{Gauss--Bonnet and the Torus}
\label{subsec:gauss-bonnet-torus}
%%%%%%%%%%%%%%%%%%%%%%%%%%%%%%%%%%%%%%%%%%%%%%%%%%%%%%%%%%%%

The torus satisfies

\[
\chi(T^2)=0.
\]

Therefore

\[
\boxed{
\int_{T^2}K\,dA=0.
}
\]

A standard embedded torus contains regions of positive and negative Gaussian
curvature.

Gauss--Bonnet requires these contributions to balance globally so that the
total signed curvature is zero.

%%%%%%%%%%%%%%%%%%%%%%%%%%%%%%%%%%%%%%%%%%%%%%%%%%%%%%%%%%%%
\subsection{Curvature and Topology}
\label{subsec:curvature-and-topology}
%%%%%%%%%%%%%%%%%%%%%%%%%%%%%%%%%%%%%%%%%%%%%%%%%%%%%%%%%%%%

Gauss--Bonnet demonstrates a central theme of modern geometry:

\[
\boxed{
\text{local geometric information}
\quad\Longrightarrow\quad
\text{global topological restrictions}.
}
\]

For a closed orientable surface of genus \(g\),

\[
\boxed{
\chi=2-2g.
}
\]

Hence

\[
\boxed{
\int_MK\,dA
=
4\pi(1-g).
}
\]

Thus the total curvature depends only on the topology of the surface.

%%%%%%%%%%%%%%%%%%%%%%%%%%%%%%%%%%%%%%%%%%%%%%%%%%%%%%%%%%%%
\subsection{Isometries}
\label{subsec:riemannian-isometries}
%%%%%%%%%%%%%%%%%%%%%%%%%%%%%%%%%%%%%%%%%%%%%%%%%%%%%%%%%%%%

\begin{definitionbox}{Riemannian Isometry}
A smooth map

\[
F:M\to N
\]

between Riemannian manifolds is an \emph{isometry} if it preserves the
Riemannian metric:

\[
\boxed{
g_M(v,w)
=
g_N(dF(v),dF(w)).
}
\]
\end{definitionbox}

An isometry preserves intrinsic geometric quantities such as:

\begin{itemize}
    \item lengths;
    \item angles;
    \item geodesic distance;
    \item Gaussian or sectional curvature;
    \item geodesic structure.
\end{itemize}

%%%%%%%%%%%%%%%%%%%%%%%%%%%%%%%%%%%%%%%%%%%%%%%%%%%%%%%%%%%%
\subsection{Local Isometries}
\label{subsec:local-isometries}
%%%%%%%%%%%%%%%%%%%%%%%%%%%%%%%%%%%%%%%%%%%%%%%%%%%%%%%%%%%%

A map may preserve the metric locally without being globally one-to-one.

For example, the plane can locally model a cylinder because a cylinder can
be unrolled without stretching.

Thus a cylinder is locally isometric to the plane.

This explains intrinsically why both have

\[
K=0.
\]

%%%%%%%%%%%%%%%%%%%%%%%%%%%%%%%%%%%%%%%%%%%%%%%%%%%%%%%%%%%%
\subsection{Riemannian Volume}
\label{subsec:riemannian-volume}
%%%%%%%%%%%%%%%%%%%%%%%%%%%%%%%%%%%%%%%%%%%%%%%%%%%%%%%%%%%%

A Riemannian metric determines a natural volume element.

In local coordinates,

\[
\boxed{
dV_g
=
\sqrt{\det(g_{ij})}\,
dx^1\cdots dx^n.
}
\]

For a two-dimensional surface this reduces to

\[
\sqrt{EG-F^2}\,du\,dv.
\]

Thus the metric determines both distance and volume.

%%%%%%%%%%%%%%%%%%%%%%%%%%%%%%%%%%%%%%%%%%%%%%%%%%%%%%%%%%%%
\subsection{Gradient on a Riemannian Manifold}
\label{subsec:riemannian-gradient}
%%%%%%%%%%%%%%%%%%%%%%%%%%%%%%%%%%%%%%%%%%%%%%%%%%%%%%%%%%%%

For a smooth function

\[
f:M\to\R,
\]

the Riemannian gradient

\[
\operatorname{grad}f
\]

is defined by

\[
\boxed{
g(\operatorname{grad}f,X)=df(X)
}
\]

for every tangent vector field \(X\).

In coordinates,

\[
\boxed{
(\operatorname{grad}f)^i
=
g^{ij}
\frac{\partial f}{\partial x^j}.
}
\]

Thus the familiar Euclidean gradient depends implicitly on the Euclidean
metric.

%%%%%%%%%%%%%%%%%%%%%%%%%%%%%%%%%%%%%%%%%%%%%%%%%%%%%%%%%%%%
\subsection{Laplace--Beltrami Operator}
\label{subsec:laplace-beltrami}
%%%%%%%%%%%%%%%%%%%%%%%%%%%%%%%%%%%%%%%%%%%%%%%%%%%%%%%%%%%%

The ordinary Laplacian can also be generalized to curved spaces.

The resulting operator is the \emph{Laplace--Beltrami operator}

\[
\boxed{
\Delta_g f.
}
\]

In local coordinates,

\[
\boxed{
\Delta_g f
=
\frac{1}{\sqrt{|g|}}
\frac{\partial}{\partial x^i}
\left(
\sqrt{|g|}\,
g^{ij}
\frac{\partial f}{\partial x^j}
\right),
}
\]

where

\[
|g|=\det(g_{ij}).
\]

This operator is central in geometric analysis and mathematical physics.

%%%%%%%%%%%%%%%%%%%%%%%%%%%%%%%%%%%%%%%%%%%%%%%%%%%%%%%%%%%%
\subsection{Differential Geometry and General Relativity}
\label{subsec:differential-geometry-relativity}
%%%%%%%%%%%%%%%%%%%%%%%%%%%%%%%%%%%%%%%%%%%%%%%%%%%%%%%%%%%%

General relativity uses a related but not positive-definite metric geometry.

Spacetime is modeled by a Lorentzian manifold.

Einstein's field equation is

\[
\boxed{
G_{\mu\nu}
+
\Lambda g_{\mu\nu}
=
\frac{8\pi G}{c^4}
T_{\mu\nu}.
}
\]

Here:

\begin{itemize}
    \item \(g_{\mu\nu}\) is the spacetime metric;
    \item \(G_{\mu\nu}\) is the Einstein tensor;
    \item \(T_{\mu\nu}\) is the stress-energy tensor;
    \item \(\Lambda\) is the cosmological constant;
    \item \(G\) is Newton's gravitational constant;
    \item \(c\) is the speed of light.
\end{itemize}

The physical theory belongs beyond the present geometry section, but it
illustrates the importance of curvature in modern science.

%%%%%%%%%%%%%%%%%%%%%%%%%%%%%%%%%%%%%%%%%%%%%%%%%%%%%%%%%%%%
\subsection{Differential Geometry and Optimization}
\label{subsec:differential-geometry-optimization}
%%%%%%%%%%%%%%%%%%%%%%%%%%%%%%%%%%%%%%%%%%%%%%%%%%%%%%%%%%%%

Many optimization problems have constraints that naturally form manifolds.

Examples include:

\[
\|\mathbf{x}\|=1,
\]

orthogonal matrices,

\[
Q^TQ=I,
\]

and positive-definite matrices.

Instead of leaving the constraint surface, one can perform optimization
directly on the manifold.

This leads to \emph{Riemannian optimization}.

A typical update uses tangent-space information followed by a geometric map
back to the manifold.

%%%%%%%%%%%%%%%%%%%%%%%%%%%%%%%%%%%%%%%%%%%%%%%%%%%%%%%%%%%%
\subsection{Differential Geometry and Mechanics}
\label{subsec:differential-geometry-mechanics}
%%%%%%%%%%%%%%%%%%%%%%%%%%%%%%%%%%%%%%%%%%%%%%%%%%%%%%%%%%%%

The possible configurations of a mechanical system often form a manifold
called its \emph{configuration space}.

Examples include:

\begin{itemize}
    \item a pendulum: \(S^1\);
    \item orientation of a rigid body: \(SO(3)\);
    \item robotic joint configurations: products of circles and other
    manifolds.
\end{itemize}

Dynamics can then be formulated geometrically on these spaces.

%%%%%%%%%%%%%%%%%%%%%%%%%%%%%%%%%%%%%%%%%%%%%%%%%%%%%%%%%%%%
\subsection{Differential Geometry and Lie Groups}
\label{subsec:differential-geometry-lie-groups}
%%%%%%%%%%%%%%%%%%%%%%%%%%%%%%%%%%%%%%%%%%%%%%%%%%%%%%%%%%%%

A Lie group is simultaneously:

\begin{itemize}
    \item a smooth manifold; and
    \item a group.
\end{itemize}

Examples include

\[
GL(n,\R),
\qquad
SO(n),
\qquad
SU(n).
\]

Differential geometry studies their tangent spaces, invariant metrics,
geodesics, and curvature.

The tangent space at the identity forms the corresponding Lie algebra.

%%%%%%%%%%%%%%%%%%%%%%%%%%%%%%%%%%%%%%%%%%%%%%%%%%%%%%%%%%%%
\subsection{Differential Geometry and Geometric Analysis}
\label{subsec:differential-geometry-geometric-analysis}
%%%%%%%%%%%%%%%%%%%%%%%%%%%%%%%%%%%%%%%%%%%%%%%%%%%%%%%%%%%%

Geometric analysis studies geometric problems using tools from analysis,
especially differential equations.

Important examples include:

\begin{itemize}
    \item minimal-surface equations;
    \item harmonic maps;
    \item mean curvature flow;
    \item Ricci flow;
    \item geometric PDEs;
    \item spectral geometry.
\end{itemize}

These subjects combine differential geometry with Partial Differential
Equations, Functional Analysis, and Calculus of Variations.

%%%%%%%%%%%%%%%%%%%%%%%%%%%%%%%%%%%%%%%%%%%%%%%%%%%%%%%%%%%%
\subsection{Mean Curvature Flow}
\label{subsec:mean-curvature-flow-preview}
%%%%%%%%%%%%%%%%%%%%%%%%%%%%%%%%%%%%%%%%%%%%%%%%%%%%%%%%%%%%

Under mean curvature flow, a hypersurface evolves approximately in the
normal direction according to its mean curvature.

Schematically,

\[
\boxed{
\frac{\partial\mathbf{X}}{\partial t}
=
-H\mathbf{n}
}
\]

under one common sign convention.

Highly curved regions move more rapidly.

This geometric evolution equation has important applications in geometry,
topology, and image processing.

%%%%%%%%%%%%%%%%%%%%%%%%%%%%%%%%%%%%%%%%%%%%%%%%%%%%%%%%%%%%
\subsection{Ricci Flow}
\label{subsec:ricci-flow-preview}
%%%%%%%%%%%%%%%%%%%%%%%%%%%%%%%%%%%%%%%%%%%%%%%%%%%%%%%%%%%%

Ricci flow evolves a Riemannian metric according to

\[
\boxed{
\frac{\partial g}{\partial t}
=
-2\operatorname{Ric}.
}
\]

The equation tends to redistribute curvature.

Ricci flow became especially famous through its role in the proof of the
Poincar\'e conjecture.

Its rigorous study belongs to Geometric Analysis.

%%%%%%%%%%%%%%%%%%%%%%%%%%%%%%%%%%%%%%%%%%%%%%%%%%%%%%%%%%%%
\subsection{Common Errors}
\label{subsec:differential-geometry-common-errors}
%%%%%%%%%%%%%%%%%%%%%%%%%%%%%%%%%%%%%%%%%%%%%%%%%%%%%%%%%%%%

\subsubsection{Confusing Speed and Curvature}

The quantity

\[
\|\gamma'(t)\|
\]

measures speed.

Curvature measures the rate at which the tangent direction changes with arc
length.

They are different geometric quantities.

%%%%%%%%%%%%%%%%%%%%%%%%%%%%%%%%%%%%%%%%%%%%%%%%%%%%%%%%%%%%
\subsubsection{Assuming Curvature Is \(f''(x)\)}

For a graph \(y=f(x)\),

\[
f''(x)
\]

alone is not geometric curvature.

The correct unsigned curvature is

\[
\boxed{
\kappa
=
\frac{|f''|}
{(1+(f')^2)^{3/2}}.
}
\]

%%%%%%%%%%%%%%%%%%%%%%%%%%%%%%%%%%%%%%%%%%%%%%%%%%%%%%%%%%%%
\subsubsection{Confusing Curve Curvature and Surface Curvature}

The curvature

\[
\kappa
\]

of a curve is different from the Gaussian curvature

\[
K
\]

of a surface.

A curve has one bending direction, while a surface has many tangent
directions and therefore several related curvature quantities.

%%%%%%%%%%%%%%%%%%%%%%%%%%%%%%%%%%%%%%%%%%%%%%%%%%%%%%%%%%%%
\subsubsection{Assuming Visible Bending Means Nonzero Gaussian Curvature}

A cylinder visibly bends in \(\R^3\), but

\[
K=0.
\]

Gaussian curvature measures intrinsic curvature rather than simply visual
extrinsic bending.

%%%%%%%%%%%%%%%%%%%%%%%%%%%%%%%%%%%%%%%%%%%%%%%%%%%%%%%%%%%%
\subsubsection{Assuming Every Geodesic Is Globally Shortest}

Geodesics are locally straight and locally minimizing under appropriate
conditions.

They need not minimize distance globally.

%%%%%%%%%%%%%%%%%%%%%%%%%%%%%%%%%%%%%%%%%%%%%%%%%%%%%%%%%%%%
\subsubsection{Treating Christoffel Symbols as Tensor Components}

Christoffel symbols

\[
\Gamma^k_{ij}
\]

do not transform as the components of a tensor.

They are coordinate-dependent coefficients of a connection.

Curvature tensors constructed from them do transform tensorially.

%%%%%%%%%%%%%%%%%%%%%%%%%%%%%%%%%%%%%%%%%%%%%%%%%%%%%%%%%%%%
\subsubsection{Assuming Zero Christoffel Symbols Mean Zero Curvature}

Christoffel symbols can vanish at a single point in normal coordinates even
on a curved manifold.

Curvature cannot generally be removed by a coordinate change.

%%%%%%%%%%%%%%%%%%%%%%%%%%%%%%%%%%%%%%%%%%%%%%%%%%%%%%%%%%%%
\subsubsection{Confusing Intrinsic and Extrinsic Geometry}

The first fundamental form describes intrinsic metric geometry.

The second fundamental form depends on how a surface is embedded and
describes extrinsic bending.

%%%%%%%%%%%%%%%%%%%%%%%%%%%%%%%%%%%%%%%%%%%%%%%%%%%%%%%%%%%%
\subsubsection{Ignoring Sign Conventions}

Several formulas in differential geometry depend on convention, especially:

\begin{itemize}
    \item the orientation of the unit normal;
    \item the sign of the shape operator;
    \item mean curvature;
    \item the Riemann curvature tensor.
\end{itemize}

A consistent convention should be maintained throughout a calculation.

%%%%%%%%%%%%%%%%%%%%%%%%%%%%%%%%%%%%%%%%%%%%%%%%%%%%%%%%%%%%
\subsection{Applications and Connections}
\label{subsec:differential-geometry-applications}
%%%%%%%%%%%%%%%%%%%%%%%%%%%%%%%%%%%%%%%%%%%%%%%%%%%%%%%%%%%%

\begin{applicationbox}
\textbf{Multivariable Calculus.}
Differential geometry gives geometric meaning to derivatives, tangent
planes, surface integrals, and coordinate changes.

\medskip

\textbf{Linear Algebra.}
Tangent spaces, metric tensors, shape operators, eigenvalues, bilinear forms,
and curvature computations rely heavily on linear algebra.

\medskip

\textbf{Non-Euclidean Geometry.}
Hyperbolic, Euclidean, and spherical geometries become standard examples of
constant-curvature Riemannian geometry.

\medskip

\textbf{Topology.}
Gauss--Bonnet connects local curvature with the global topology of surfaces.

\medskip

\textbf{Differential Topology.}
Smooth manifolds, tangent spaces, differential forms, and smooth maps provide
the common language of the two subjects.

\medskip

\textbf{Partial Differential Equations.}
Minimal surfaces, harmonic functions, curvature flows, and geometric
evolution equations lead naturally to PDEs.

\medskip

\textbf{Calculus of Variations.}
Geodesics and minimal surfaces arise as critical points of geometric
functionals.

\medskip

\textbf{Mathematical Physics.}
Riemannian and Lorentzian geometry provide mathematical frameworks for
mechanics, field theory, and relativity.

\medskip

\textbf{Optimization.}
Riemannian optimization extends gradient-based methods to nonlinear
constraint manifolds.

\medskip

\textbf{Robotics and Aerospace.}
Rotations and rigid-body configurations naturally live on manifolds such as
\(SO(3)\).

\medskip

\textbf{Computer Graphics.}
Surface normals, curvature, meshes, geodesics, and shape analysis are central
to geometric modeling.

\medskip

\textbf{Data Science.}
Manifold learning and geometric data analysis attempt to exploit nonlinear
low-dimensional structures within high-dimensional data.
\end{applicationbox}

%%%%%%%%%%%%%%%%%%%%%%%%%%%%%%%%%%%%%%%%%%%%%%%%%%%%%%%%%%%%
\subsection{Exercises}
\label{subsec:differential-geometry-exercises}
%%%%%%%%%%%%%%%%%%%%%%%%%%%%%%%%%%%%%%%%%%%%%%%%%%%%%%%%%%%%

\begin{exercise}[1]
For

\[
\gamma(t)=(t,t^2),
\]

compute

\[
\gamma'(t).
\]
\end{exercise}

\begin{exercise}[1]
Determine whether

\[
\gamma(t)=(t^2,t^3)
\]

is regular at \(t=0\).
\end{exercise}

\begin{exercise}[1]
Find the speed of

\[
\gamma(t)=(3\cos t,3\sin t).
\]
\end{exercise}

\begin{exercise}[1]
Find the arc length of

\[
\gamma(t)=(3\cos t,3\sin t),
\qquad
0\leq t\leq2\pi.
\]
\end{exercise}

\begin{exercise}[2]
Compute the unit tangent vector of

\[
\gamma(t)=(t,t^2)
\]

at \(t=1\).
\end{exercise}

\begin{exercise}[2]
Compute the curvature of

\[
y=x^2
\]

at \(x=1\).
\end{exercise}

\begin{exercise}[2]
Find the curvature of

\[
y=x^3
\]

at \(x=0\).
\end{exercise}

\begin{exercise}[2]
Show that a circle of radius \(R\) has constant curvature

\[
\kappa=\frac1R.
\]
\end{exercise}

\begin{exercise}[2]
Find the radius of curvature of a circle of radius \(5\).
\end{exercise}

\begin{exercise}[2]
For

\[
\gamma(t)=(\cos t,\sin t,t),
\]

compute

\[
\gamma'(t),
\qquad
\gamma''(t).
\]
\end{exercise}

\begin{exercise}[3]
Compute the curvature and torsion of

\[
\gamma(t)=(\cos t,\sin t,t).
\]
\end{exercise}

\begin{exercise}[2]
Explain geometrically why a plane curve has torsion zero wherever the
Frenet frame is defined.
\end{exercise}

\begin{exercise}[2]
For the surface

\[
\mathbf{X}(u,v)=(u,v,u^2+v^2),
\]

compute

\[
\mathbf{X}_u
\qquad\text{and}\qquad
\mathbf{X}_v.
\]
\end{exercise}

\begin{exercise}[2]
For the previous surface, compute

\[
\mathbf{X}_u\times\mathbf{X}_v.
\]
\end{exercise}

\begin{exercise}[2]
Find a unit normal to

\[
z=x^2+y^2
\]

at the origin.
\end{exercise}

\begin{exercise}[2]
For the cylinder

\[
\mathbf{X}(\theta,z)
=
(R\cos\theta,R\sin\theta,z),
\]

compute \(E,F,G\).
\end{exercise}

\begin{exercise}[2]
Use the previous exercise to show that

\[
ds^2=R^2\,d\theta^2+dz^2.
\]
\end{exercise}

\begin{exercise}[3]
Show that the Gaussian curvature of a cylinder is zero.
\end{exercise}

\begin{exercise}[2]
For a sphere of radius \(R\), verify that

\[
dA
=
R^2\sin\phi\,d\theta\,d\phi.
\]
\end{exercise}

\begin{exercise}[2]
Integrate the spherical area element to recover

\[
A=4\pi R^2.
\]
\end{exercise}

\begin{exercise}[2]
A surface has principal curvatures

\[
\kappa_1=2,
\qquad
\kappa_2=3.
\]

Find its Gaussian and mean curvatures.
\end{exercise}

\begin{exercise}[2]
A surface has

\[
\kappa_1=4,
\qquad
\kappa_2=-2.
\]

Find \(K\) and \(H\), and classify the point according to the sign of \(K\).
\end{exercise}

\begin{exercise}[2]
A surface has principal curvatures

\[
\frac12
\qquad\text{and}\qquad
0.
\]

Find its Gaussian curvature.
\end{exercise}

\begin{exercise}[3]
Explain why rolling a plane into a cylinder does not change Gaussian
curvature.
\end{exercise}

\begin{exercise}[3]
Explain why no open region of a sphere can be flattened onto the plane by a
distance-preserving map.
\end{exercise}

\begin{exercise}[2]
What are the geodesics of the Euclidean plane?
\end{exercise}

\begin{exercise}[2]
What are the geodesics of a sphere?
\end{exercise}

\begin{exercise}[3]
Explain why a great-circle arc need not be globally shortest between its
endpoints.
\end{exercise}

\begin{exercise}[2]
For the Euclidean metric

\[
g=
\begin{pmatrix}
1&0\\
0&1
\end{pmatrix},
\]

compute all Christoffel symbols in Cartesian coordinates.
\end{exercise}

\begin{exercise}[3]
Use the result of the previous exercise to derive the Euclidean geodesic
equation.
\end{exercise}

\begin{exercise}[2]
Explain why Christoffel symbols can vanish at a point even when curvature
does not.
\end{exercise}

\begin{exercise}[2]
State the difference between a tangent space and a cotangent space.
\end{exercise}

\begin{exercise}[2]
Expand the Einstein-summation expression

\[
g_{ij}v^iw^j
\]

explicitly when \(i,j\in\{1,2\}\).
\end{exercise}

\begin{exercise}[2]
If

\[
g=
\begin{pmatrix}
4&0\\
0&9
\end{pmatrix}
\]

and

\[
v=
\begin{pmatrix}
1\\2
\end{pmatrix},
\]

compute

\[
g(v,v).
\]
\end{exercise}

\begin{exercise}[2]
Use the previous result to find the metric length of \(v\).
\end{exercise}

\begin{exercise}[3]
Explain conceptually why a Riemannian metric is a smoothly varying family of
inner products.
\end{exercise}

\begin{exercise}[3]
Describe what parallel transport around a closed loop can reveal about
curvature.
\end{exercise}

\begin{exercise}[2]
Compute

\[
\chi(S^2).
\]

Then use Gauss--Bonnet to determine the total Gaussian curvature of any
Riemannian metric on \(S^2\).
\end{exercise}

\begin{exercise}[2]
Use

\[
\chi(T^2)=0
\]

to determine the total Gaussian curvature of a torus.
\end{exercise}

\begin{exercise}[2]
A closed orientable surface has genus \(g=2\).

Compute its Euler characteristic and total Gaussian curvature.
\end{exercise}

\begin{exercise}[3]
For a closed orientable surface of genus \(g\), show that

\[
\int_MK\,dA
=
4\pi(1-g).
\]
\end{exercise}

\begin{exercise}[3]
Explain why a closed orientable surface of genus \(g\geq2\) cannot have
Gaussian curvature strictly positive everywhere.
\end{exercise}

\begin{exercise}[3]
Explain how the exponential map converts tangent vectors into points reached
by geodesics.
\end{exercise}

\begin{exercise}[3]
What does geodesic completeness mean?
Give one complete and one incomplete example.
\end{exercise}

\begin{exercise}[3]
Explain conceptually how Jacobi fields measure the separation of nearby
geodesics.
\end{exercise}

\begin{exercise}[3]
Why do geodesics tend to diverge more strongly in negatively curved spaces?
\end{exercise}

\begin{exercise}[3]
Explain the difference among Gaussian curvature, sectional curvature, Ricci
curvature, and scalar curvature.
\end{exercise}

\begin{exercise}[3]
Explain why a minimal surface need not have zero Gaussian curvature.
\end{exercise}

\begin{exercise}[3]
Show that a sphere cannot be a minimal surface in Euclidean three-space.
\end{exercise}

\begin{exercise}[3]
Explain how geodesics can be viewed as critical points of the energy
functional.
\end{exercise}

\begin{exercise}[3]
Explain why differential forms are useful for integration on manifolds.
\end{exercise}

\begin{exercise}[3]
Explain how generalized Stokes' theorem extends the fundamental theorem of
calculus.
\end{exercise}

\begin{exercise}[4]
Derive the Christoffel symbols for the polar-coordinate metric

\[
ds^2=dr^2+r^2d\theta^2.
\]
\end{exercise}

\begin{exercise}[4]
Use the Christoffel symbols from the previous exercise to write the geodesic
equations in polar coordinates.
\end{exercise}

\begin{exercise}[4]
Investigate the Gaussian curvature of the graph

\[
z=x^2-y^2.
\]

Explain why the origin is a hyperbolic point.
\end{exercise}

\begin{exercise}[4]
Investigate the principal curvatures of a sphere and show that every tangent
direction is a principal direction.
\end{exercise}

\begin{exercise}[4]
Research the relationship among Ricci curvature, volume growth, and geodesic
behavior.
\end{exercise}

\begin{exercise}[4]
Research how Ricci flow was used in the proof of the Poincar\'e conjecture.
\end{exercise}

%%%%%%%%%%%%%%%%%%%%%%%%%%%%%%%%%%%%%%%%%%%%%%%%%%%%%%%%%%%%
\subsection{Conceptual Summary}
\label{subsec:differential-geometry-summary}
%%%%%%%%%%%%%%%%%%%%%%%%%%%%%%%%%%%%%%%%%%%%%%%%%%%%%%%%%%%%

Differential geometry applies calculus and linear algebra to curves,
surfaces, and manifolds.

For a regular curve,

\[
\boxed{
\mathbf{T}
=
\frac{\gamma'}{\|\gamma'\|}
}
\]

is the unit tangent vector.

Its curvature is

\[
\boxed{
\kappa
=
\left\|
\frac{d\mathbf{T}}{ds}
\right\|.
}
\]

For space curves, the Frenet frame consists of

\[
\boxed{
\mathbf{T},
\qquad
\mathbf{N},
\qquad
\mathbf{B},
}
\]

while torsion

\[
\boxed{\tau}
\]

measures twisting out of the osculating plane.

For a parametrized surface

\[
\mathbf{X}(u,v),
\]

the tangent plane is generated by

\[
\boxed{
\mathbf{X}_u,
\qquad
\mathbf{X}_v.
}
\]

The first fundamental form is

\[
\boxed{
I
=
E\,du^2+2F\,du\,dv+G\,dv^2
}
\]

and determines intrinsic metric geometry.

The second fundamental form is

\[
\boxed{
II
=
e\,du^2+2f\,du\,dv+g\,dv^2
}
\]

and describes extrinsic bending.

The principal curvatures are

\[
\boxed{
\kappa_1,\qquad\kappa_2.
}
\]

Gaussian and mean curvature are

\[
\boxed{
K=\kappa_1\kappa_2,
}
\]

\[
\boxed{
H=\frac{\kappa_1+\kappa_2}{2}.
}
\]

Gauss's Theorema Egregium shows that Gaussian curvature is intrinsic.

A Riemannian manifold carries a smoothly varying inner product

\[
\boxed{
g_p:T_pM\times T_pM\to\R.
}
\]

The Levi-Civita connection allows tangent vectors to be differentiated
intrinsically.

Geodesics satisfy

\[
\boxed{
\nabla_{\dot{\gamma}}\dot{\gamma}=0,
}
\]

or, in coordinates,

\[
\boxed{
\frac{d^2x^k}{dt^2}
+
\Gamma^k_{ij}
\frac{dx^i}{dt}
\frac{dx^j}{dt}
=
0.
}
\]

Curvature in higher dimensions is encoded by the Riemann curvature tensor,
from which sectional, Ricci, and scalar curvature are obtained.

Finally, Gauss--Bonnet gives the profound relationship

\[
\boxed{
\int_MK\,dA
=
2\pi\chi(M),
}
\]

linking local geometry with global topology.

Differential geometry therefore forms a major bridge among calculus, linear
algebra, topology, differential equations, analysis, and mathematical
physics.

%%%%%%%%%%%%%%%%%%%%%%%%%%%%%%%%%%%%%%%%%%%%%%%%%%%%%%%%%%%%
\subsection{Section Connections}
\label{subsec:differential-geometry-connections}
%%%%%%%%%%%%%%%%%%%%%%%%%%%%%%%%%%%%%%%%%%%%%%%%%%%%%%%%%%%%

\chapterconnection{
Differential Geometry extends the study of Euclidean and Non-Euclidean
Geometry by allowing curvature and metric structure to vary smoothly from
point to point. Multivariable and Vector Calculus provide the local
derivative and integration machinery, while Linear Algebra provides tangent
spaces, bilinear forms, eigenvalues, and tensor operations. The manifold
viewpoint leads naturally to Differential Topology and Manifold Theory.
Gauss--Bonnet provides a fundamental connection with Algebraic and Geometric
Topology. Riemannian metrics, geodesics, curvature tensors, and differential
operators later reappear in Geometric Analysis, Partial Differential
Equations, Calculus of Variations, Lie Theory, Mathematical Physics, and
General Relativity. The next section, Algebraic Geometry, studies geometric
spaces using polynomial equations and algebraic structures rather than
primarily differential methods.
}

%%%%%%%%%%%%%%%%%%%%%%%%%%%%%%%%%%%%%%%%%%%%%%%%%%%%%%%%%%%%
% SECTION SOURCE TRACKING
%%%%%%%%%%%%%%%%%%%%%%%%%%%%%%%%%%%%%%%%%%%%%%%%%%%%%%%%%%%%

%------------------------------------------------------------
% 3.7 Differential Geometry
%------------------------------------------------------------
%
% Source basis:
%   - Standard differential geometry and introductory
%     Riemannian geometry.
%
% Used for:
%   - regular parametrized curves;
%   - arc length and unit-speed parametrization;
%   - tangent, normal, and binormal vectors;
%   - curvature and torsion;
%   - Frenet--Serret equations;
%   - regular parametrized surfaces;
%   - tangent planes and normals;
%   - first and second fundamental forms;
%   - principal, Gaussian, and mean curvature;
%   - shape operator;
%   - intrinsic and extrinsic geometry;
%   - Theorema Egregium;
%   - minimal surfaces;
%   - geodesics and energy;
%   - smooth manifolds, charts, and tangent spaces;
%   - Riemannian metrics;
%   - Levi-Civita connection and Christoffel symbols;
%   - geodesic equation;
%   - parallel transport;
%   - Riemann, sectional, Ricci, and scalar curvature;
%   - exponential map and completeness;
%   - Hopf--Rinow theorem;
%   - Jacobi fields;
%   - differential forms and generalized Stokes theorem;
%   - Gauss--Bonnet theorem;
%   - geometric-analysis connections.
%
% Notes:
%   - Multivariable differentiation, vector calculus, and
%     surface integration are developed in Analysis.
%   - The formal construction of smooth manifolds is expanded
%     in Manifold Theory and Differential Topology.
%   - Differential forms and de Rham theory receive fuller
%     treatment in Differential Topology and Algebraic Topology.
%   - Tensor analysis and curvature are introduced here
%     geometrically; advanced tensor calculus belongs to later
%     geometric-analysis and applied-mathematics material.
%   - Lorentzian geometry and Einstein's equations are included
%     only as connections to Mathematical Physics.
%   - Minimal surfaces, Ricci flow, mean curvature flow, and
%     geometric PDEs are developed more systematically in
%     Geometric Analysis and PDEs.
%
%%%%%%%%%%%%%%%%%%%%%%%%%%%%%%%%%%%%%%%%%%%%%%%%%%%%%%%%%%%%